BFV 支持在整数明文模数下对明文多项式进行精确运算。与 CKKS 的近似运算不同，只要每一步均满足 BFV 正确性界，其解码结果就是精确的；仅仅采用整数编码，并不能保证任意噪声或电路深度下都能正确解密。应用还必须选择能够保持预期整数语义的明文模数与编码。本章以 Fan 和 Vercauteren 的原始方案~\cite{cryptoeprint:2012/144} 为参考。

在 BFV 中，每个明文都被加密为一个 RLWE 密文。因此，BFV 的密文与密文加法、密文与明文加法，以及密文与明文乘法，均基于 GLWE 的同态加法与乘法（见 \autoref{part:generic-fhe}）实现；取 $k = 1$ 即可将 GLWE 化为 RLWE。




\begin{tcolorbox}[
    title = \textbf{所需背景知识},    % box title
    colback = white,    % light background; tweak to taste
    colframe = black,  % frame colour
    boxrule = 0.8pt,     % line thickness
    left = 1mm, right = 1mm, top = 1mm, bottom = 1mm % inner padding
]

\begin{itemize}
\item \autoref{sec:modulo}: \nameref{sec:modulo}
\item \autoref{sec:group}: \nameref{sec:group}
\item \autoref{sec:field}: \nameref{sec:field}
\item \autoref{sec:order}: \nameref{sec:order}
\item \autoref{sec:polynomial-ring}: \nameref{sec:polynomial-ring}
\item \autoref{sec:decomp}: \nameref{sec:decomp}
\item \autoref{sec:roots}: \nameref{sec:roots}
\item \autoref{sec:cyclotomic}: \nameref{sec:cyclotomic}
\item \autoref{sec:cyclotomic-polynomial-integer-ring}: \nameref{sec:cyclotomic-polynomial-integer-ring}
\item \autoref{sec:matrix}: \nameref{sec:matrix}
\item \autoref{sec:euler}: \nameref{sec:euler}
\item \autoref{sec:modulus-rescaling}: \nameref{sec:modulus-rescaling}
\item \autoref{sec:chinese-remainder}: \nameref{sec:chinese-remainder}
\item \autoref{sec:polynomial-interpolation}: \nameref{sec:polynomial-interpolation}
\item \autoref{sec:ntt}: \nameref{sec:ntt}
\item \autoref{sec:lattice}: \nameref{sec:lattice}
\item \autoref{sec:rlwe}: \nameref{sec:rlwe}
\item \autoref{sec:glwe}: \nameref{sec:glwe}
\item \autoref{sec:glwe-add-cipher}: \nameref{sec:glwe-add-cipher}
\item \autoref{sec:glwe-add-plain}: \nameref{sec:glwe-add-plain}
\item \autoref{sec:glwe-mult-plain}: \nameref{sec:glwe-mult-plain}
\item \autoref{subsec:modulus-switch-rlwe}: \nameref{subsec:modulus-switch-rlwe}
\item \autoref{sec:glwe-key-switching}: \nameref{sec:glwe-key-switching}
\end{itemize}
\end{tcolorbox}

\clearpage

\subsection{单值编码}
\label{subsec:bfv-single-encoding}

这种二进制多项式编码要求 $0\leq m<2^n$。只有选定的整数系数代表没有模 $t$ 回绕，且乘积次数严格小于 $n$ 时，在 $X=2$ 处求值才保持整数加法和乘法。它不是从 $\mathbb Z_t[X]/(X^n+1)$ 到无界整数的环同态。例如，取 $n=2$ 时，整数 $2$ 的编码 $X$ 在商环中平方得到 $-1$，而不是整数 $4$。

BFV 支持两种编码方案：单值编码和批量编码。本小节介绍单值编码。



\begin{tcolorbox}[title={\textbf{\tboxlabel{\ref*{subsec:bfv-single-encoding}} BFV 编码}}]

\textbf{\underline{输入整数}：}将满足 $0\leq m<2^n$ 的整数 $m$ 分解如下：

$m = b_{n-1}\cdot 2^{n-1} + b_{n-2}\cdot 2^{n-2} + \cdots + b_1\cdot 2^1 + b_0\cdot 2^0$  \text { } ，其中每个 $b_i \in \{0, 1\}$

$ $


\textbf{\underline{编码后的多项式}：} $M(X) = b_0 + b_1X + b_2X^2 + \cdots + b_{n-1}X^{n-1} \in \mathcal{R}_{\langle n, t\rangle}$

$ $

\textbf{\underline{解码}：} $M(X=2) = m$

\end{tcolorbox}


下面分析小结~\ref*{subsec:bfv-single-encoding} 中的编码方案，是否能保证加法和乘法之后仍能正确解码。这等价于证明：

$\text{Decode}(\sigma(m_1) + \sigma(m_2)) = m_1 + m_2$ \textcolor{red}{$\rhd$ 其中 $\sigma$ 为编码函数}

$\text{Decode}(\sigma(m_1) \cdot \sigma(m_2)) = m_1 \cdot m_2$

设 $\sigma(m_1) = M_1(X)$ 且 $\sigma(m_2) = M_2(X)$，则

$\text{Decode}(\sigma(m_1) + \sigma(m_2)) = \text{Decode}(M_1(X) + M_2(X))$


$= \text{Decode}(M_{1+2}(X))$ \textcolor{red}{$\rhd$ 其中 $M_{1+2}(X) = M_1(X) + M_2(X)$}

$= M_{1+2}(2)$ \textcolor{red}{$\rhd$ 因为解码就是在 $X = 2$ 处对多项式求值}

$= M_{1}(2) + M_{2}(2)$ \textcolor{red}{$\rhd$ 因为在 $X=2$ 处对 $M_{1+2}(X)$ 求值，等价于将 $M_{1+2}(X)$ 拆分为 $M_{1}(X)$ 和 $M_{2}(X)$，分别对 $M_1(2)$ 与 $M_2(2)$ 求值后相加}

$ $

类似地，乘法之后的解码也保持正确：

$\text{Decode}(\sigma(m_1) \cdot \sigma(m_2)) = \text{Decode}(M_1(X) \cdot M_2(X))$

$= \text{Decode}(M_{1\cdot2}(X))$ \textcolor{red}{$\rhd$ 其中 $M_{1\cdot2}(X) = M_1(X) \cdot M_2(X)$}

$= M_{1\cdot2}(2)$ \textcolor{red}{$\rhd$ 因为解码就是在 $X = 2$ 处对多项式求值}

$= M_{1}(2) \cdot M_{2}(2)$ \textcolor{red}{$\rhd$ 因为在 $X=2$ 处对 $M_{1\cdot2}(X)$ 求值，等价于将 $M_{1\cdot2}(X)$ 拆分为 $M_{1}(X)$ 和 $M_{2}(X)$，分别对 $M_1(2)$ 与 $M_2(2)$ 求值后相乘}




$ $

因此，这些整数恒等式在所述无回绕及次数条件下成立，并不是商环中无条件成立的恒等式。

$ $

总结~\ref*{subsec:bfv-single-encoding} 的编码只有在所有结果多项式的次数严格小于 $n$，且所选整数系数不发生模 $t$ 回绕时，才保持整数语义。次数达到 $n$ 时，模 $X^n+1$ 约减就已经改变在 $2$ 处的求值。
由于这些约束，单值编码并不是全同态加密的理想选择。而且，尽管每个多项式包含 $n$ 个系数，却只能编码一个值，因此计算效率也较低。


\subsection{批量编码}
\label{subsec:bfv-batch-encoding}

下述标量槽位公式假设 $n\geq2$ 是二次幂，且 $t$ 是满足 $2n\mid(t-1)$ 的素数。此时 $\omega=g^{(t-1)/(2n)}$ 是本原 $2n$ 次单位根，范德蒙矩阵可逆。若 $t=p^r$ 是素数幂，且 $p\equiv1\pmod{2n}$，将根提升到 $\mathbb Z_{p^r}$ 后也能得到这些标量槽位；直接使用 $t-1$ 的指数公式一般不适用于素数幂。其他参数可能只能得到数量较少的扩环槽位。本章的 $W^T$ 是按行排列求值点的范德蒙矩阵，$W$ 是其转置；这里显式规定矩阵朝向。

单值编码方案（\autoref{subsec:bfv-single-encoding}）每次只对一个值进行编码与解码，而批量编码利用高维向量，同时对大量数值执行这些操作。因此，批量编码比单值编码更高效。此外，通过向量加法和 Hadamard 积，可以同时对批量编码的值逐元素执行同态加法与乘法。因此，批量编码能够采用 SIMD（单指令多数据）方式，比单值编码更高效地完成同态运算。

BFV 编码将模 $t$ 的、含 $n$ 个槽位的整数输入向量 $\vec{v} = (v_0, v_1, v_2, \cdots v_{n-1})$，转换为模 $t$ 的另一个含 $n$ 个槽位的向量 $\vec{m} = (m_0, m_1, m_2, \cdots m_{n-1})$。后者的各元素作为编码所得 $(n-1)$ 次或更低次数多项式 $M(X) \in \mathbb{Z}_t[X] / (X^n + 1)$ 的系数。

$ $

\subsubsection{编码第一阶段（\textsf{Encoding\textsubscript{1}}）}
\label{subsubsec:bfv-encoding-1} 

在 \autoref{subsec:poly-vector-transformation} 中，我们已经知道，通过映射 $\sigma$，可以将 $(n-1)$ 次或更低次数的多项式同构地映射为一个含 $n$ 个槽位的向量（为简化记号，将 \autoref{subsec:poly-vector-transformation} 中的 $\sigma_c$ 记为 $\sigma$）：

$\sigma: M(X) \in \mathbb{Z}_t[X]/(X^n + 1) \longrightarrow  (M(\omega),M(\omega^3),M(\omega^5), \cdots, M(\omega^{2n-1})) \in \mathbb{Z}_t^n$

$ $

，即在 $n$ 个不同的 $(\mu=2n)$ 次本原单位根 $\omega, \omega^3, \omega^5, \cdots, \omega^{2n-1}$ 处对多项式 $M(X)$ 求值。令 $\vec{m}$ 为包含多项式 $M(X)$ 的 $n$ 个系数的向量，则可将映射 $\sigma$ 表示为：

$\vec{v} = W^T \cdot \vec{m}$

$ $

，其中 $W^T$ 如下：

$W^T = \begin{bmatrix}
1 & (\omega) & (\omega)^2 & \cdots & (\omega)^{n-1}\\
1 & (\omega^3) & (\omega^3)^2 & \cdots & (\omega^3)^{n-1}\\
1 & (\omega^5) & (\omega^5)^2 & \cdots & (\omega^5)^{n-1}\\
\vdots & \vdots & \vdots & \ddots & \vdots \\
1 & (\omega^{2n-1}) & (\omega^{2n-1})^2 & \cdots & (\omega^{2n-1})^{n-1}\\
\end{bmatrix}$  \textcolor{red}{ \text{ } \# $W^T$ 是 \autoref{subsec:poly-vector-transformation} 中所述 $W$ 的转置}

$ $

注意，$W^T$ 的每一行与 $\vec{m}$ 的点积，给出多项式 $M(X)$ 在相应 $X = \{w, w^3, w^5, \cdots, w^{2n-1}\}$ 处的求值结果。在 BFV 编码方案中，\textsf{Encoding\textsubscript{1}} 将含 $n$ 个槽位的输入向量 $\vec{v}$ $\in \mathbb{Z}_t^n$ 编码为明文多项式 $M(X) \in \mathbb{Z}_t[X] / (X^n + 1)$，而 \textsf{Decoding\textsubscript{2}} 则将 $M(X)$ 解码回 $\vec{v}$。由于 $W^T \cdot \vec{m}$ 得到的 $\vec{v}$ 是 $M(X)$ 的解码，因此将 $W^T$ 称为解码矩阵。另一方面，\textsf{Encoding\textsubscript{1}} 的目标是将 $\vec{v}$ 编码为 $M(X)$，以便基于 $M(X)$ 进行同态计算。由关系 $\vec{v} = W^T \cdot \vec{m}$，可以推导出编码公式：

$(W^T)^{-1} \cdot \vec{v} = (W^T)^{-1} W^T \cdot \vec{m}$

$\vec{m} = (W^T)^{-1} \cdot \vec{v}$

$ $

因此，需要求出 $W^T$ 的逆矩阵 $(W^T)^{-1}$，作为编码矩阵。根据 \autoref{subsec:vandermonde-euler} 中的定理~\ref*{subsec:vandermonde-euler}，已经知道 $V^{-1} = \dfrac{V^T \cdot I_n^R}{n}$，其中 $V = W^T$ 且 $V^T = W$。换言之，$(W^T)^{-1} = \dfrac{W \cdot I_n^R}{n}$。因此，BFV 的编码公式可以写为：

$\vec{m} = (W^T)^{-1} \cdot \vec{v} = \dfrac{W \cdot I_n^R \cdot \vec{v}}{n}$，\text{ } 其中

$ $

$W =  \begin{bmatrix}
1 & 1 & 1 & \cdots & 1\\
(\omega) & (\omega^3) & (\omega^5) & \cdots & (\omega^{2n-1})\\
(\omega)^2 & (\omega^3)^2 & (\omega^5)^2 & \cdots & (\omega^{2n-1})^2\\
\vdots & \vdots & \vdots & \ddots & \vdots \\
(\omega)^{n-1} & (\omega^3)^{n-1} & (\omega^5)^{n-1} & \cdots & (\omega^{2n-1})^{n-1}\\
\end{bmatrix}$

$= \begin{bmatrix}
1 & 1 & \cdots & 1 & 1 & \cdots & 1 & 1\\
(\omega) & (\omega^3) & \cdots & (\omega^{n - 1}) & (\omega^{-(n - 1)}) & \cdots & (\omega^{-3}) & (\omega^{-1})\\
(\omega)^2 & (\omega^3)^2 & \cdots & (\omega^{n - 1})^2 & (\omega^{-(n - 1)})^2 & \cdots & (\omega^{-3})^2 & (\omega^{-1})^2\\
\vdots & \vdots & \vdots & \ddots & \vdots \\
(\omega)^{n-1} & (\omega^3)^{n-1} & \cdots & (\omega^{n - 1})^{n-1} & (\omega^{-(n - 1)})^{n-1} & \cdots & (\omega^{-3})^{n-1} & (\omega^{-1})^{n-1}\\
\end{bmatrix}$

$ $

，其中 $\omega$ 是模 $t$ 的 $2n$ 次本原单位根（这意味着 $t \equiv 1 \bmod 2n$）。因此有 $\omega = g^{\frac{t - 1}{2n}} \bmod t$（$g$ 是 $\mathbb{Z}_t^{\times}$ 的生成元，见 \autoref{subsubsec:poly-vector-transformation-modulus}）。

在 \autoref{subsec:poly-vector-transformation} 中，我们已经知道 $W$ 是 $n$ 维向量空间的一组基。因此，对于 $n$ 维向量空间 $\mathbb{Z}_t^{n}$ 中的每个 $\vec{v}$，都存在唯一对应的向量 $\dfrac{W \cdot I_n^R\cdot\vec{v}}{n} = \vec{m}$（参见 \autoref{subsec:projection} 中的定理~\ref*{subsec:projection}）。由此，以 $\vec{m}$ 的 $n$ 个元素为系数构成的多项式 $M(X)$，也是与 $\vec{v}$ 双射对应的唯一多项式。

注意，通过计算 $\dfrac{W \cdot I_n^R \cdot \vec{v}}{n}$
，这给出线性双射。该环同构把 $\vec v$ 的逐槽乘法对应到 $M$ 的负循环多项式乘法，而不是系数向量 $\vec m$ 的逐分量乘法。

\subsubsection{编码第二阶段（\textsf{Encoding\textsubscript{2}}）}

得到含 $n$ 个槽位的向量 $\vec{m}$ 后，使用尺度因子 $\Delta = \left\lfloor \dfrac{q}{t} \right\rfloor$ 对其放大（即相乘），其中 $q$ 是密文模数。将 $\vec{m}$ 乘以 $\Delta$，得到 $\Delta \vec{m}$。
$\Delta\vec{m}$ 中的 $n$ 个整数将作为 RLWE 加密中明文多项式的 $n$ 个系数。最终编码后的明文多项式为 $\Delta M = \sum\limits_{i=0}^{n-1}\Delta m_iX^i$。


\subsubsection{解码第一阶段（\textsf{Decoding\textsubscript{1}}）}
\label{subsubsec:bfv-enc-dec-decoding1} 

编码与解码示意图中的除以 $\Delta$，针对的是理想、未约减且无噪声的编码多项式。对于实际密文，应先使用 \autoref{subsec:bfv-enc-dec} 的缩放舍入解码器，再对 $\mathcal R_{\langle n,t\rangle}$ 中恢复出的多项式应用槽位求值矩阵。

RLWE 密文完成解密的前半部分、得到 $\Delta M = \sum\limits_{i=0}^{n-1}\Delta m_iX^i$ 后，计算 $\dfrac{\Delta\vec{m}}{\Delta} = \vec{m}$。




\subsubsection{解码第二阶段（\textsf{Decoding\textsubscript{2}}）}


在 \autoref{subsubsec:bfv-encoding-1} 中，我们已经推导出以下解码公式，用于将一个系数为模 $t$ 整数、次数为 $(n-1)$ 的多项式转换成包含 $n$ 个槽位的输入向量：

$\vec{v} = W^T \cdot \vec{m}$


\subsubsection{小结}
\label{subsubsec:bfv-encoding-summary}

\begin{tcolorbox}[title={\textbf{\tboxlabel{\ref*{subsubsec:bfv-encoding-summary}} BFV 的编码与解码}}]



\textbf{\underline{输入}：} 包含 $n$ 个槽位的模 $t$ 整数向量 $\vec{v} = (v_0, v_1, \cdots, v_{n-1}) \in \mathbb{Z}_t^n$

\par\noindent\rule{\textwidth}{0.4pt}

\textbf{\underline{编码}}
\begin{enumerate}

\item 应用变换 $\vec{m} = n^{-1}\cdot W \cdot I_n^R \cdot \vec{v}$，将 $\vec{v} \in \mathbb{Z}_t^n$ 转换为 $\vec{m} \in \mathbb{Z}_t^n$

，其中，$W$ 是按如下方式构造的 $n$ 维向量空间的一组基：

$W =  \begin{bmatrix}
1 & 1 & 1 & \cdots & 1\\
(\omega) & (\omega^3) & (\omega^5) & \cdots & (\omega^{2n-1})\\
(\omega)^2 & (\omega^3)^2 & (\omega^5)^2 & \cdots & (\omega^{2n-1})^2\\
\vdots & \vdots & \vdots & \ddots & \vdots \\
(\omega)^{n-1} & (\omega^3)^{n-1} & (\omega^5)^{n-1} & \cdots & (\omega^{2n-1})^{n-1}\\
\end{bmatrix}$

$= \begin{bmatrix}
1 & 1 & \cdots & 1 & 1 & \cdots & 1 & 1\\
(\omega) & (\omega^3) & \cdots & (\omega^{n - 1}) & (\omega^{-(n - 1)}) & \cdots & (\omega^{-3}) & (\omega^{-1})\\
(\omega)^2 & (\omega^3)^2 & \cdots & (\omega^{n - 1})^2 & (\omega^{-(n - 1)})^2 & \cdots & (\omega^{-3})^2 & (\omega^{-1})^2\\
\vdots & \vdots & \vdots & \ddots & \vdots \\
(\omega)^{n-1} & (\omega^3)^{n-1} & \cdots & (\omega^{n - 1})^{n-1} & (\omega^{-(n - 1)})^{n-1} & \cdots & (\omega^{-3})^{n-1} & (\omega^{-1})^{n-1}\\
\end{bmatrix}$

$ $

，其中，$\omega$ 是模 $t$ 的一个本原 $2n$ 次单位根。这意味着 $\omega = g^{\frac{t - 1}{2n}} \bmod t$（$g$ 是 $\mathbb{Z}_t^{\times}$ 的一个生成元，见 \autoref{subsubsec:poly-vector-transformation-modulus}）。

$ $

\item 将 $\vec{m}$ 转换为缩放后的整数向量 $\Delta\vec{m}$，其中，$\Delta=\lfloor \dfrac{q}{t}\rfloor$ 是尺度因子。如果 $\Delta$ 过于接近 1（即 $t$ 太大），则\textbf{噪声预算}会过小（容易导致解密失败）。如果 $\Delta$ 过于接近 $q$（即 $t$ 太小），则\textbf{消息容量}会过小（即明文模数 $t$ 限制了可编码数值的范围）。最终编码得到的明文多项式为 $\Delta M = \sum\limits_{i=0}^{n-1} \Delta m_i X^i \text{ } \in \mathbb{Z}_q[X] / (X^n + 1)$。 %The rounding of $\lceil \Delta \vec{m} \rfloor$ during the encoding process makes CKKS an approximate FHE scheme.

\end{enumerate}

\par\noindent\rule{\textwidth}{0.4pt}

\textbf{\underline{解码}：} 从明文多项式 $\Delta M = \sum\limits_{i=0}^{n-1}\Delta m_iX^i$ 中恢复 $\vec{m} = \dfrac{\Delta \vec{m}}{\Delta}$。然后，
计算 $\vec{v} = W^T \cdot \vec{m}$。

\end{tcolorbox}


%\para{BFV's Encoding Range:} \autoref{tab:polynomial-encoding-capacity-integer} depicts the range of input values that can be encoded by the polynomials $M(X) \in \mathcal{R}_{\langle n, t \rangle}$ over $X \in \mathbb{Z}_t$. As illustrated in \autoref{tab:polynomial-encoding-capacity-integer}, the encoding range of integer inputs is determined by the size of the $n$ and $t$ parameters. Therefore, these two parameters should be chosen sufficiently large enough to encode all values that can be used by an application. 

不过，小结~\ref*{subsubsec:bfv-encoding-summary} 还不是 BFV 批量编码的最终版本。在 \autoref{subsec:bfv-rotation} 中，我们将介绍如何在不解密封装输入向量的密文的情况下，对其槽位进行同态旋转。为了支持这种同态旋转，需要对小结~\ref*{subsubsec:bfv-encoding-summary} 中介绍的编码方案稍作调整。我们将在 \autoref{subsec:bfv-rotation} 中说明具体方法，并在 \autoref{subsubsec:bfv-rotation-summary} 的小结~\ref*{subsubsec:bfv-rotation-summary} 中给出 BFV 的最终编码方案。


\subsection{加密与解密}
\label{subsec:bfv-enc-dec}

这里使用 RLWE 相位 $B+AS$ 描述对称 BFV 加密，即 \autoref{subsec:glwe-alternative} 中取 $k=1$ 的另一种符号约定（\autoref{sec:rlwe}）。用 $[U]_q$ 表示逐系数中心模约减，并设 $r_q=q-t\Delta$。完整解码器 $\textsf{Dec}$ 使用比例 $t/q$。本章其他相位推导中的 $\textsf{RLWE}^{-1}$ 表示最终缩放舍入之前的带噪相位。任何一种记号都不表示同态运算后的噪声仍是新鲜噪声。



\begin{tcolorbox}[title={\textbf{\tboxlabel{\ref*{subsec:bfv-enc-dec}} BFV 加密与解密}}]

\textbf{\underline{初始设置}：} $\Delta=\left\lfloor\dfrac{q}{t}\right\rfloor \text{ 是多项式编码的明文尺度因子}, \text{ } S \xleftarrow{\$} \mathcal{R}_{\langle n, \textit{tern} \rangle}$

基本方案的明文模数可以是任意满足 $q\gg t$ 的整数 $t\geq2$。批编码和素数幂自举方法还需满足各自小节所列的附加条件。这里 $S$ 的系数取自三元集合 $\{-1,0,1\}$。

\par\noindent\rule{\textwidth}{0.4pt}

\textbf{\underline{加密输入}：} $\Delta M \in \mathcal{R}_{\langle n, q \rangle}$、$A \xleftarrow{\$} \mathcal{R}_{\langle n, q \rangle}$、$E \xleftarrow{\chi_\sigma} \mathcal{R}_{\langle n, q \rangle}$


$ $

\begin{enumerate}


\item 计算 $B = -A \cdot S + \Delta M + E \text{ } \in \mathcal{R}_{\langle n,q \rangle}$

\item $\textsf{RLWE}_{S,\sigma}(\Delta M + E) = (A, B) \text{ } \in \mathcal{R}_{\langle n,q \rangle}^2$ 

\end{enumerate}

\par\noindent\rule{\textwidth}{0.4pt}

\textbf{\underline{解密输入}：} $\textsf{ct} = (A, B) \text{ } \in \mathcal{R}_{\langle n,q \rangle}^2$

$\textsf{Dec}_S(A,B)=\left[\left\lceil\dfrac{t}{q}[B+AS]_q\right\rfloor\right]_t=M$。

（噪声 $E = \sum\limits_{i=0}^{n-1}e_iX^i$ 由舍入过程消除。）



$ $

\textbf{\underline{正确解密的条件}：}

对于相位 $\Delta M+E\pmod q$，一个逐系数充分条件是 $\|tE-r_qM\|_\infty<q/2$。特别地，当 $M$ 使用满足 $\|M\|_\infty\leq t/2$ 的中心代表，且右侧为正时，$\|E\|_\infty<q/(2t)-r_q/2$ 是充分条件。同态运算后对明文模约减还会产生余数误差，推导见 \autoref{subsubsec:bfv-noise-analysis}。这是 BFV 的缩放舍入条件；不能在不改变假设的情况下，用 \autoref{subsubsec:lwe-noise-bound} 讨论的另一种 $1/\Delta$ 约定替代它。

\end{tcolorbox}


在本节中，为简洁起见，我们经常将 $\textsf{RLWE}_{S,\sigma}(\Delta  M + E)$ 写作 $\textsf{RLWE}_{S,\sigma}(\Delta  M)$，因为 $\textsf{RLWE}_{S,\sigma}(\Delta M + E) \approx \textsf{RLWE}_{S,\sigma}(\Delta M)$（即二者解密得到相同的消息）。即使我们用 $\textsf{RLWE}_{S,\sigma}(\Delta  M)$ 代替 $\textsf{RLWE}_{S,\sigma}(\Delta  M + E)$，也应将其理解为对 $\Delta  M + E$ 的加密（即噪声包含在缩放后的消息中）。


我们将在 \autoref{subsubsec:bfv-noise-analysis} 中更详细地说明 BFV 正确解密的条件。



\subsection{密文与密文相加}
\label{subsec:bfv-add-cipher}

BFV 的密文与密文相加采用 RLWE 的密文与密文相加方案，但将加密和解密公式中 $A\cdot S$ 项的符号取反。具体而言，这等价于 $k = 1$ 时的另一种 GLWE 版本（\autoref{subsec:glwe-alternative}）的密文与密文相加方案。


\begin{tcolorbox}[title={\textbf{\tboxlabel{\ref*{subsec:bfv-add-cipher}} BFV 密文与密文相加}}]
$\textsf{RLWE}_{S, \sigma}(\Delta M^{\langle 1 \rangle} + E^{\langle 1 \rangle} ) + \textsf{RLWE}_{S, \sigma}(\Delta M^{\langle 2 \rangle} + E^{\langle 2 \rangle}) $

$ = ( A^{\langle 1 \rangle}, \text{ } B^{\langle 1 \rangle}) + (A^{\langle 2 \rangle}, \text{ } B^{\langle 2 \rangle}) $

$ = ( A^{\langle 1 \rangle} + A^{\langle 2 \rangle}, \text{ } B^{\langle 1 \rangle} + B^{\langle 2 \rangle} ) $

$= \textsf{RLWE}_{S, \sigma}(\Delta(M^{\langle 1 \rangle} + M^{\langle 2 \rangle}) + E^{\langle 1 \rangle} + E^{\langle 2 \rangle})$
\end{tcolorbox}


\subsubsection{噪声界分析}
\label{subsubsec:bfv-noise-analysis}

令 $r_q=q-t\Delta$，中心相位为 $x=[\Delta M+E]_q=\Delta M+E-qK$。乘以 $t/q$ 得到精确恒等式
\[
\frac{t}{q}x=M+\frac{tE-r_qM}{q}-tK.
\]
最后模 $t$ 约减时，整数倍 $tK$ 消失。因此，\autoref{subsec:bfv-enc-dec} 中的总结~\ref*{subsec:bfv-enc-dec} 在以下条件下正确解密：
\[
\|tE-r_qM\|_\infty<q/2.
\]
该证明没有删去内部模 $q$ 约减。当 $r_q\ne0$ 时，用 $1/\Delta$ 代替 $t/q$ 并不能得到相同恒等式，而需要另行指定解码器和代表元约定。

对于密文加法，写成 $M^{\langle1\rangle}+M^{\langle2\rangle}=M^{\langle1+2\rangle}+tH$，其中 $M^{\langle1+2\rangle}$ 是所选中心明文代表。因为 $t\Delta=q-r_q$，相加密文的相位为
\[
\Delta M^{\langle1+2\rangle}+E_{\rm add}\pmod q,
\qquad E_{\rm add}=E^{\langle1\rangle}+E^{\langle2\rangle}-r_qH.
\]
其正确性的充分条件为
\[
\bigl\|t(E^{\langle1\rangle}+E^{\langle2\rangle}-r_qH)-r_qM^{\langle1+2\rangle}\bigr\|_\infty<q/2.
\]
项 $-r_qH$ 记录明文回绕，不能省略。乘法或旋转后，只要明确有效明文和完整误差，也可应用相同推理。这与 \autoref{subsubsec:lwe-noise-bound} 中需要显式保留余数项的原则一致。增大 $q/t$ 可能增加正确性预算，但安全性仍需根据完整参数集评估。

\subsection{密文与明文相加}
\label{subsec:bfv-add-plain}

相加的明文多项式使用尺度 $\Delta$。若 $M+\Lambda=M_{\rm out}+tH$，则明文约减后的有效误差为 $E-r_qH$，其中 $r_q=q-t\Delta$，与 \autoref{subsubsec:bfv-noise-analysis} 中的约定相同。

BFV 的密文与明文相加采用 RLWE 的密文与明文相加方案，但将加密和解密公式中 $A\cdot S$ 项的符号取反。具体而言，这等价于 $k = 1$ 时的另一种 GLWE 版本（\autoref{subsec:glwe-alternative}）的密文与明文相加方案（\autoref{sec:glwe-add-plain}）。


\begin{tcolorbox}[title={\textbf{\tboxlabel{\ref*{subsec:bfv-add-plain}} BFV 密文与明文相加}}]
$\textsf{RLWE}_{S, \sigma}(\Delta M + E) + \Delta\Lambda $

$=  (A, \text{ } B) + \Delta\Lambda$

$=  (A, \text{ } B + \Delta\cdot\Lambda)$

$= \textsf{RLWE}_{S, \sigma}(\Delta (M + \Lambda) + E)$
\end{tcolorbox}



\subsection{密文与明文相乘}
\label{subsec:bfv-mult-plain}

这里 $\Lambda$ 是未缩放的明文乘数。若 $M\Lambda=M_{\rm out}+tH$，则有效误差为 $\Lambda E-r_qH$。对于多项式乘法，$\|\Lambda E\|_\infty\leq\|\Lambda\|_1\|E\|_\infty$；这个卷积界及余数项都必须计入解码条件。


BFV 的密文与明文相乘采用 RLWE 的密文与明文相乘方案，但将加密和解密公式中 $A\cdot S$ 项的符号取反。具体而言，这等价于 $k = 1$ 时的另一种 GLWE 版本（\autoref{subsec:glwe-alternative}）的密文与明文相乘方案（\autoref{sec:glwe-mult-plain}）。


\begin{tcolorbox}[title={\textbf{\tboxlabel{\ref*{subsec:bfv-mult-plain}} BFV 密文与明文相乘}}]
$\textsf{RLWE}_{S, \sigma}(\Delta M + E) \cdot \Lambda$

$= (A, \text{ } B) \cdot \Lambda$

$= (A \cdot \Lambda, \text{ }  B \cdot \Lambda )$

$= \textsf{RLWE}_{S, \sigma}(\Delta (M \cdot \Lambda) + \Lambda E)$
\end{tcolorbox}




\subsection{密文与密文相乘}
\label{subsec:bfv-mult-cipher}

以下四阶段表述是一个需要显式误差界的教学变体：它使用辅助模数 $Q=q\Delta$，并以 $1/\Delta$ 重缩放。原始 FV 乘法在重线性化之前，以 $t/q$ 缩放整数提升后的张量乘积分量；两种比例只有在 $t\mid q$ 时才相同。这里保留该变体以解释模提升，但完整保留其附加余数和舍入误差，并不声称所有 BFV 实现都按此顺序执行这四个阶段。

\noindent \textbf{- 参考资料 1：}
\href{https://www.inferati.com/blog/fhe-schemes-bfv}{BFV 加密方案简介}~\cite{inferati-bfv}

\noindent \textbf{- 参考资料 2：}
\href{https://eprint.iacr.org/2012/144.pdf}{某种程度上、部分与全同态加密}~\cite{cryptoeprint:2012/144}

%\noindent \textbf{- Reference 3:} 
%\href{https://eprint.iacr.org/2016/510.pdf}{A Full RNS Variant $of FV like Somewhat Homomorphic Encryption Schemes}~\cite{cryptoeprint:2016/510}


给定两个密文 $\textsf{RLWE}_{S, \sigma}(\Delta M^{\langle 1\rangle })$ 和 $\textsf{RLWE}_{S, \sigma}(\Delta M^{\langle 2\rangle})$，密文与密文相乘的目标是构造一个新密文，使其解密结果为 $\Delta M^{\langle 1 \rangle} M^{\langle 2 \rangle}$。密文与密文相乘比密文与明文相乘更复杂，包括四个步骤：（1）\textsf{ModRaise}；（2）多项式乘法；（3）重线性化；（4）重缩放。

为便于理解，我们将基于 RLWE 的另一种版本（\autoref{subsec:glwe-alternative} 中的定理~\ref*{subsec:glwe-alternative}）来解释 BFV 的密文与密文相乘；在该版本中，加密和解密公式中 $AS$ 项的符号取反。

\subsubsection{模数提升（\textsf{ModRaise}）}
\label{subsubsec:bfv-mult-cipher-modraise}

选定 $(A,B)$ 的整数系数代表，在 $\mathbb Z[X]/(X^n+1)$ 中将未约减相位写为
\[
AS+B=\Delta M+E+qK.
\]
把相同系数代表嵌入模数 $Q=q\Delta$ 后，相位变为 $\Delta M+E+qK\pmod Q$，并不是只有误差 $E$ 的、原明文 $\Delta M$ 的加密。一般不能断言 $(q-1)n+(q-1)<Q$；这还要求 $n$ 与 $\Delta$ 满足额外关系。上述同余本身也不需要无回绕条件。后续两个阶段的乘积和相位都按模 $Q$ 解释；重缩放时会显式计入整数提升及其中的 $Q$ 倍数。

BFV 密文与密文相乘的第一步是对两个输入密文 $(A^{\langle 1 \rangle}, B^{\langle 1 \rangle}) \bmod q$ 和 $(A^{\langle 2 \rangle}, B^{\langle 2 \rangle}) \bmod q$ 进行\textit{模数提升}，
按如下方式将模数变为 $q \rightarrow Q$（其中 $Q = q\cdot \Delta$）：

$(A^{\langle 1 \rangle}, B^{\langle 1 \rangle}) \bmod Q$

$(A^{\langle 2 \rangle}, B^{\langle 2 \rangle}) \bmod Q$

$ $

经过 \textsf{ModRaise} 后，这两个密文的解密结果为：

$A^{\langle 1 \rangle} \cdot S + B^{\langle 1 \rangle} = \Delta M^{\langle 1 \rangle} + E^{\langle 1 \rangle} + K_1q \pmod Q$

$A^{\langle 2 \rangle} \cdot S + B^{\langle 2 \rangle} = \Delta M^{\langle 2 \rangle} + E^{\langle 2 \rangle} + K_2q \pmod Q$


$ $

因此，\textit{模数提升后}的密文具有以下形式：

$\textsf{RLWE}_{S, \sigma}(\Delta M^{\langle 1 \rangle} + K_1q) = (A^{\langle 1 \rangle}, B^{\langle 1 \rangle}) \bmod Q$

$\textsf{RLWE}_{S, \sigma}(\Delta M^{\langle 2 \rangle} + K_2q) = (A^{\langle 2 \rangle}, B^{\langle 2 \rangle}) \bmod Q$


\subsubsection{多项式乘法}
\label{subsubsec:bfv-mult-cipher-multiplication}

接下来的目标是构造一个新密文，它加密的内容为
$(\Delta M^{\langle 1 \rangle} + E^{\langle 1 \rangle} + K_1q) \cdot (\Delta M^{\langle 2 \rangle} + E^{\langle 2 \rangle} + K_2q)$. 

首先，可以推导出以下关系：

$(\Delta M^{\langle 1 \rangle} + E^{\langle 1 \rangle} + K_1q) \cdot (\Delta M^{\langle 2 \rangle} + E^{\langle 2 \rangle} + K_2q)$

$ = (A^{\langle 1 \rangle} \cdot S + B^{\langle 1 \rangle}) \cdot (A^{\langle 2 \rangle} \cdot S + B^{\langle 2 \rangle})$

$ = \underbrace{B^{\langle 1 \rangle}B^{\langle 2 \rangle}}_{D_0}  + \underbrace{(B^{\langle 2 \rangle}A^{\langle 1 \rangle} + B^{\langle 1 \rangle}A^{\langle 2 \rangle})}_{D_1} \cdot S + \underbrace{(A^{\langle 1 \rangle} \cdot A^{\langle 2 \rangle})}_{D_2} \cdot S \cdot S $



$= D_0 + D_1\cdot S + D_2\cdot S^2$



$ $

同时，我们还具有以下关系：

$\textsf{RLWE}_{S, \sigma}^{-1}(\Delta M^{\langle 1 \rangle} + K_1q) = \Delta M^{\langle 1 \rangle} + E^{\langle 1 \rangle} + K_1q$

$\textsf{RLWE}_{S, \sigma}^{-1}(\Delta M^{\langle 2 \rangle} + K_2q) = \Delta M^{\langle 2 \rangle} + E^{\langle 2 \rangle} + K_2q$

$ $

将这些关系结合起来，得到：

$\textsf{RLWE}_{S, \sigma}^{-1}(\Delta M^{\langle 1 \rangle} + K_1q) \cdot \textsf{RLWE}_{S, \sigma}^{-1}(\Delta M^{\langle 2 \rangle} + K_2q) = D_0 + D_1\cdot S + D_2\cdot S^2$

$ $

注意，作为密文分量的 $D_0, D_1,$ 和 $D_2$ 是已知值，而 $S$ 只有私钥持有者知道。因此，可以将 $D_0 + D_1\cdot S + D_2\cdot S^2$ 视为解密公式：给定密文分量 $D_0, D_1, D_2$ 和私钥 $S$，便可得到 $(\Delta M^{\langle 1 \rangle} + E^{\langle 1 \rangle} + K_1q) \cdot (\Delta M^{\langle 2 \rangle} + E^{\langle 2 \rangle} + K_2q)$。换言之，可以令 $(D_0, D_1, D_2)$ 为一种新的密文形式，它能够由 $S$ 解密为 $(\Delta M^{\langle 1 \rangle} + E^{\langle 1 \rangle} + K_1q) \cdot (\Delta M^{\langle 2 \rangle} + E^{\langle 2 \rangle} + K_2q)$。

三元组 $(D_0,D_1,D_2)$ 是关于秘密密钥 $S$ 的二次密文，因此解密使用的是平方 $S^2$，不是平方根。继续相乘会使秘密密钥次数相加；重复平方则在每步将次数翻倍。重线性化以引入求值密钥误差为代价，将密文恢复为两个分量及一次形式。



\subsubsection{重线性化}
\label{subsubsec:bfv-mult-cipher-relinearization}

对于任意 $Q$，本小节使用整数分解基 $g_i=\beta^{i-1}$，取 $l=\lceil\log_\beta Q\rceil$，逐系数精确分解满足 $D_2\equiv\sum_{i=1}^lD_{2,i}g_i\pmod Q$。这是递增权重约定，不是递减的 $Q/\beta^i$ 约定；相匹配的 RLev 各行加密 $g_iS^2$。令 $E_{\rm rl}=\sum_iD_{2,i}E_i'$ 表示累计误差。重线性化后的相位等于原相位加 $E_{\rm rl}$，并非与原相位完全相等。

重线性化将三元组 $(D_0,D_1,D_2)\in\mathcal R_{\langle n,Q\rangle}^3$ 转换为两个分量。取 $\textsf{ct}_\alpha=(D_1,D_0)$，并通过分解构造 $\textsf{ct}_\beta$，则有
\[
\phi_S(\textsf{ct}_\alpha+\textsf{ct}_\beta)\equiv D_0+D_1S+D_2S^2+E_{\rm rl}\pmod Q.
\]

在 $D_0+D_1S+D_2S^2$ 中，前两项构成合成密文 $\textsf{ct}_\alpha=(D_1,D_0)$。分解构造产生的 $\textsf{ct}_\beta$ 满足相位关系 $\phi_S(\textsf{ct}_\beta)\equiv D_2S^2+E_{\rm rl}\pmod Q$。公开秘密密钥函数（如 $S^2$）的加密，还需要所选构造的求值密钥安全性假设；下列代数并不构成对该假设的证明。

构造加密 $D_2 \cdot S^2$ 的密文有一种直接的方法：将 $S^2$ 加密为 RLWE 密文 $\textsf{RLWE}_{S, \sigma}(S^2) = (A^{\langle s \rangle}, B^{\langle s \rangle})$，使其满足 $A^{\langle s \rangle}\cdot S + B^{\langle s \rangle} = S^2 + E^{\langle s \rangle} \bmod Q$（其中，密文模数为 $Q$，明文尺度因子为 $\Delta = 1$）。然后，将 $D_2$ 看作模 $Q$ 的明文多项式，与 $D_2$ 进行密文与明文相乘（\autoref{sec:glwe-mult-plain}）。但这种方法在实际中不可行，因为计算 $D_2 \cdot \textsf{RLWE}_{S, \sigma}(S^2)$ 会产生如下巨大的噪声：

$D_2 \cdot (A^{\langle s \rangle}, B^{\langle s \rangle}) = (D_2\cdot A^{\langle s \rangle}, D_2 \cdot B^{\langle s \rangle})$

$ $

，其解密结果为：

$D_2\cdot A^{\langle s \rangle} \cdot S + D_2\cdot B^{\langle s \rangle} = D_2 \cdot S^2 + D_2\cdot E^{\langle s \rangle} \pmod Q$ 

$ $

。在上述解密表达式 $D_2S^2 + D_2E^{\langle s \rangle} \bmod Q$ 中，$D_2S^2$ 项可以模 $Q$ 约减，因为在最终解密公式 $D_0 + D_1S + D_2S^2 \bmod Q$ 中，这一项本来就允许模 $Q$ 约减。然而，问题出在噪声 $D_2\cdot E^{\langle s \rangle}$：其系数可以取 $[0, Q - 1]$ 中的任意值（因为多项式 $D_2 = A^{\langle 1 \rangle} A^{\langle 2 \rangle}$ 的每个系数都可以取 $[0, Q - 1]$ 中的任意值）。如此大的噪声无法保证最终的正确解密。

为避免这一噪声问题，改进的方法是使用Gadget 分解技术（\autoref{subsec:gadget-decomposition}），将加密 $D_2 S^2$ 的 RLWE 密文表示为多个噪声较小的 RLWE 密文之和。为此，我们使用加密 $S^2$ 的 RLev 密文（\autoref{sec:glev}）。假设Gadget 向量为 $\vec{g} = \Bigg(\beta^{0}, \beta^{1}, \beta^{2}, \cdots, \beta^{l-1}\Bigg)$。回顾我们的目标：在已知 $D_2$、未知 $S$ 且已知 $\textsf{RLev}_{S, \sigma}^{\beta, l}(S^2) = \left\{\textsf{RLWE}_{S, \sigma}\left(\beta^{i-1}\cdot S^2\right)\right\}_{i=1}^{l}$ 的情况下，求得 $\textsf{ct}_\beta = \textsf{RLWE}_{S, \sigma}( S^2 \cdot D_2)$。于是，可以按如下方式构造 $\textsf{ct}_\beta$：

$\textsf{ct}_\beta = \textsf{RLWE}_{S, \sigma}(S^2 \cdot D_2)$

$= \textsf{RLWE}_{S, \sigma}\left( S^2 \cdot \left(D_{2,1}\beta^{0} + D_{2,2}\beta^{1} + \cdots D_{2,l}\beta^{l-1}\right)\right)$ \textcolor{red}{ $\rhd$ 对 $D_2$ 进行分解}

$= \textsf{RLWE}_{S, \sigma}\left( S^2 \cdot D_{2,1} \cdot \beta^{0}\right) + \textsf{RLWE}_{S, \sigma}\left( S^2 \cdot D_{2,2} \cdot \beta^{1}\right) + \cdots + \textsf{RLWE}_{S, \sigma}\left( S^2 \cdot D_{2,l} \cdot \beta^{l-1}\right)$

$= D_{2,1}\cdot\textsf{RLWE}_{S, \sigma}\left( S^2 \cdot \beta^{0}\right) + D_{2,2}\cdot\textsf{RLWE}_{S, \sigma}\left( S^2 \cdot \beta^{1}\right) + \cdots + D_{2,l}\cdot\textsf{RLWE}_{S, \sigma}\left( S^2 \cdot \beta^{l-1}\right)$
\textcolor{red}{ $\rhd$ 其中，每个 \textsf{RLWE} 密文都以 $S^2\beta^{0}, S^2\beta^{1}, \cdots, S^2\beta^{l-1}$ 为明文，明文尺度因子为 $\Delta = 1$}

$ $

$= \bm{\langle} \textsf{Decomp}^{\beta, l}(D_2), \text{ } \textsf{RLev}_{S, \sigma}^{\beta, l}( S^2) \bm{\rangle}$ \textcolor{red}{ $\rhd$ 将 \textsf{Decomp} 和 \textsf{RLev} 视为向量并求内积}

$ $

对上述结果解密，得到：

$\textsf{RLWE}_{S, \sigma}^{-1}(\textsf{ct}_\beta = \bm{\langle} \textsf{Decomp}^{\beta, l}(D_2), \text{ } \textsf{RLev}_{S, \sigma}^{\beta, l}( S^2) \bm{\rangle} \bm{)}$ \textcolor{red}{ $\rhd$ $\textsf{RLev}_{S, \sigma}^{\beta, l}( S^2)$ 的尺度因子均为 1}

$= D_{2,1}\cdot\left(E_1' +  S^2\beta^{0}\right) + D_{2,2}\cdot\left(E_2' +  S^2\beta^{1}\right) + \cdots + D_{2,l}\cdot\left(E_l' +  S^2\beta^{l-1}\right)$ \textcolor{red}{ \text{ } \# 其中，每个 $E_i'$ 都是嵌入 $\textsf{RLWE}_{S, \sigma}\left(S^2\cdot\beta^{i-1}\right)$ 中的噪声}

$ $

$= \sum\limits_{i=1}^{l} (E_i'\cdot D_{2,i}) +  S^2\cdot\left(D_{2,1}\beta^{0} + D_{2,2}\beta^{1} + \cdots + D_{2, l}\beta^{l-1}\right)$

$= \sum\limits_{i=1}^{l} \epsilon_{i} + D_2\cdot S^2$ \textcolor{red}{ \text{ } \# 其中，每个 $\epsilon_i = E_i'\cdot D_{2,i}$}

$\equiv D_2S^2+E_{\rm rl}\pmod Q$，其中 $\|E_{\rm rl}\|_\infty\leq\sum_{i=1}^l\|D_{2,i}\|_1\|E_i'\|_\infty$。小数字给出有用的界，但不能与每一个可能的单密钥噪声实例作严格大小比较。

$ $

因此，我们得到以下完整关系：


$\textsf{RLWE}_{S, \sigma}^{-1}(\Delta M^{\langle 1 \rangle} + K_1q) \cdot \textsf{RLWE}_{S, \sigma}^{-1}(\Delta M^{\langle 2 \rangle} + K_2q) \bmod Q$

$= (\Delta M^{\langle 1 \rangle} + E^{\langle 1 \rangle} + K_1q) \cdot (\Delta M^{\langle 2 \rangle} + E^{\langle 2 \rangle} + K_2q) \bmod Q$

$ = (A^{\langle 1 \rangle} \cdot S + B^{\langle 1 \rangle}) \cdot (A^{\langle 2 \rangle} \cdot S + B^{\langle 2 \rangle})  \bmod Q$

$ = D_0 + D_1\cdot S + D_2\cdot S^2  \bmod Q$ \textcolor{red}{ $\rhd$ $D_0 = B^{\langle 1 \rangle} B^{\langle 2 \rangle}$, \text{ } $D_1 = A^{\langle 1 \rangle} B^{\langle 2 \rangle} + A^{\langle 2 \rangle} B^{\langle 1 \rangle}$, \text{ } $D_2 = A^{\langle 1 \rangle} A^{\langle 2 \rangle}$}

$ = \textsf{RLWE}_{S, \sigma}^{-1}(\textsf{ct}_\alpha) + \textsf{RLWE}_{S, \sigma}^{-1}(\textsf{ct}_\beta) - \sum\limits_{i=1}^{l} (E_i'\cdot D_{2,i})  \bmod Q$ 

\textcolor{red}{ $\rhd$ $\textsf{ct}_\alpha = (D_1, D_0) = (A_\alpha, B_\alpha)$, \text{ } $\textsf{ct}_\beta = \langle \textsf{Decomp}^{\beta, l}(D_2), \textsf{RLev}_{S, \sigma}^{\beta, l}(S^2) \rangle = (A_\beta, B_\beta)$}

$ $

$ = \textsf{RLWE}_{S, \sigma}^{-1}(\textsf{ct}_\alpha + \textsf{ct}_\beta) - \sum\limits_{i=1}^{l} (E_i'\cdot D_{2,i})  \bmod Q$

$ = \textsf{RLWE}_{S, \sigma}^{-1}\bm((A_{\alpha+\beta}, B_{\alpha+\beta})\bm) - \sum\limits_{i=1}^{l} (E_i'\cdot D_{2,i}) \bmod Q$  \textcolor{red}{ $\rhd$ $A_{\alpha+\beta} = A_{\alpha} + A_\beta$, \text{ } $B_{\alpha + \beta} = B_\alpha + B_\beta$}


$ $

从上述推导中提取以下关系：

$(\Delta M^{\langle 1 \rangle} + E^{\langle 1 \rangle} + K_1q) \cdot (\Delta M^{\langle 2 \rangle} + E^{\langle 2 \rangle} + K_2q) \bmod Q$

$=  \Delta^2M^{\langle 1 \rangle}M^{\langle 2 \rangle} + \Delta\cdot (M^{\langle 1 \rangle}E^{\langle 2 \rangle} + M^{\langle 2 \rangle}E^{\langle 1 \rangle}) + q\cdot(\Delta M^{\langle 1 \rangle}K_2 + \Delta M^{\langle 2 \rangle}K_1 + E^{\langle 1 \rangle}K_2 + E^{\langle 2 \rangle}K_1) + K_1 K_2 q^2 + E^{\langle 1 \rangle} E^{\langle 2 \rangle} \bmod Q$

$= \textsf{RLWE}_{S, \sigma}^{-1}\bm((A_{\alpha+\beta}, B_{\alpha+\beta})\bm) - \sum\limits_{i=1}^{l} (E_i'\cdot D_{2,i}) \bmod Q$  

$ $

可以将上述关系改写为：


$\textsf{RLWE}_{S, \sigma}^{-1}\bm((A_{\alpha+\beta}, B_{\alpha+\beta})\bm) = A_{\alpha+\beta}\cdot S + B_{\alpha+\beta} \bmod Q$

$ = \Delta^2M^{\langle 1 \rangle}M^{\langle 2 \rangle} + \Delta\cdot (M^{\langle 1 \rangle}E^{\langle 2 \rangle} + M^{\langle 2 \rangle}E^{\langle 1 \rangle}) + q\cdot(\Delta M^{\langle 1 \rangle}K_2 + \Delta M^{\langle 2 \rangle}K_1 + E^{\langle 1 \rangle}K_2 + E^{\langle 2 \rangle}K_1) + K_1 K_2 q^2 + E^{\langle 1 \rangle} E^{\langle 2 \rangle} + \sum\limits_{i=1}^{l} (E_i'\cdot D_{2,i}) \bmod Q$ 

$ $

该表达式包含乘积明文、提升项 $qK_i$、输入误差乘积及重线性化误差。下一小节会跟踪这些项经过教学变体的 $1/\Delta$ 重缩放后的变化，并应用完整 BFV 解码条件，而不预先假设它与无噪声加密近似相等。


\subsubsection{重缩放}
\label{subsubsec:bfv-mult-cipher-rescaling}

令 $(A_*,B_*)=(A_{\alpha+\beta},B_{\alpha+\beta})$ 为模数 $Q=q\Delta$ 下重线性化后的密文，并选定整数系数代表。定义
\[
C_A=\Bigl\lceil A_*/\Delta\Bigr\rfloor,\quad C_B=\Bigl\lceil B_*/\Delta\Bigr\rfloor,
\quad R_A=C_A-A_*/\Delta,\quad R_B=C_B-B_*/\Delta.
\]
$R_A,R_B$ 的每个系数绝对值至多为 $1/2$。用 $L_i=\Delta M^{\langle i\rangle}+E^{\langle i\rangle}+qK_i$ 表示整数输入相位，并保留 \autoref{subsubsec:bfv-mult-cipher-relinearization} 中的 $E_{\rm rl}$，则
\[
A_*S+B_*=L_1L_2+E_{\rm rl}+QK_*,
\qquad C_AS+C_B\equiv\frac{L_1L_2+E_{\rm rl}}{\Delta}+R_AS+R_B\pmod q.
\]
这里使用 $Q/\Delta=q$，所以被约减的 $Q$ 倍数在除法后成为整数倍 $q$。设 $r_q=q-t\Delta$，展开乘积并去除整数倍 $q$ 后得到
\[
C_AS+C_B\equiv\Delta M^{\langle1\rangle}M^{\langle2\rangle}+\eta\pmod q,
\]
\[
\begin{split}
\eta={}&M^{\langle1\rangle}E^{\langle2\rangle}+M^{\langle2\rangle}E^{\langle1\rangle}
+\frac q\Delta(E^{\langle1\rangle}K_2+E^{\langle2\rangle}K_1)\\
&+\frac{q r_q}{\Delta}K_1K_2
+\frac{E^{\langle1\rangle}E^{\langle2\rangle}+E_{\rm rl}}{\Delta}+R_AS+R_B.
\end{split}
\]
特别地，$q^2/\Delta=qt+q r_q/\Delta$。后一项一般不可忽略，不能将其视为近似零的误差而删去。$\eta$ 的表达式可以含有有理数加项，但完整表达式在模 $q$ 下与整数输出相位一致。

解码时，将 $M^{\langle1\rangle}M^{\langle2\rangle}=M_{\rm out}+tH$ 写成所选中心明文代表。有效输出误差为 $E_{\rm out}=\eta-r_qH$。标准 BFV 解码器正确性的一个充分条件是
\[
\|t(\eta-r_qH)-r_qM_{\rm out}\|_\infty<q/2.
\]
当 $t\mid q$ 时，所有 $r_q$ 项消失，但明文与噪声的乘积、求值密钥误差和各分量舍入误差仍然存在。

\para{噪声分析：}上述乘积都是负循环多项式乘积。应使用 $\|UV\|_\infty\leq\|U\|_1\|V\|_\infty\leq n\|U\|_\infty\|V\|_\infty$，不能套用逐系数标量界。另有 $\|R_AS+R_B\|_\infty\leq(\|S\|_1+1)/2$。用这些不等式约束 $\eta$ 的其余各项，再应用三角不等式与最终阈值比较。仅凭 $q\gg t$ 或 $Q=q\Delta$，不能无条件保证乘法成功。

\para{模数切换与重缩放：}在这一特定构造中，以 $1/\Delta=q/Q$ 缩放分量，会将辅助模数从 $Q$ 改回 $q$，同时降低明文尺度。两个术语描述的是不同目的，并不总能互换。标准 FV 乘法可以通过对未约减整数乘积执行缩放舍入，在恢复明文尺度的同时保持对外密文模数不变。

\subsubsection{小结}
\label{subsubsec:bfv-mult-cipher-summary}

将以上内容综合起来，BFV 的密文与密文相乘可总结如下：

\begin{tcolorbox}[title={\textbf{\tboxlabel{\ref*{subsubsec:bfv-mult-cipher-summary}} BFV 的密文与密文相乘}}]


假设有以下两个 RLWE 密文：

$\textsf{RLWE}_{S, \sigma}(\Delta M^{\langle 1 \rangle}) = (A^{\langle 1 \rangle}, B^{\langle 1 \rangle}) \bmod q$，\text{ } 其中 $B^{\langle 1 \rangle} = -A^{\langle 1 \rangle} \cdot S + \Delta M^{\langle 1 \rangle} + E^{\langle 1 \rangle} \bmod q$

$\textsf{RLWE}_{S, \sigma}(\Delta M^{\langle 2 \rangle}) = (A^{\langle 2 \rangle}, B^{\langle 2 \rangle})  \bmod q$，\text{ } 其中 $B^{\langle 2 \rangle} = -A^{\langle 2 \rangle} \cdot S + \Delta M^{\langle 2 \rangle} + E^{\langle 2 \rangle} \bmod q$

$ $

这两个密文的乘法按如下步骤执行：

$ $

\begin{enumerate}

\item \textbf{\underline{模数提升（\textsf{ModRaise}）}}

按如下方式将密文 $(A^{\langle 1 \rangle}, B^{\langle 1 \rangle}) \bmod q$ 和 $(A^{\langle 2 \rangle}, B^{\langle 2 \rangle}) \bmod q$ 的模数强制提升至 $Q$（其中 $Q = q \cdot \Delta$）：

$(A^{\langle 1 \rangle}, B^{\langle 1 \rangle}) \bmod Q$

$(A^{\langle 2 \rangle}, B^{\langle 2 \rangle}) \bmod Q$


$ $

\item \textbf{\underline{乘法}}

在模 $Q$ 下计算以下多项式乘法：

$D_0 = B^{\langle 1 \rangle}B^{\langle 2 \rangle} \bmod Q$

$D_1 = B^{\langle 2 \rangle}A^{\langle 1 \rangle} +  B^{\langle 1 \rangle}A^{\langle 2 \rangle} \bmod Q$

$D_2 = A^{\langle 1 \rangle} \cdot A^{\langle 2 \rangle} \bmod Q$


$ $

\item \textbf{\underline{重线性化}}


计算：

$ \textsf{ct}_\alpha = (D
_1, D_0)$

$ \textsf{ct}_\beta = \bm{\langle} \textsf{Decomp}^{\beta, l}(D_2), \text{ } \textsf{RLev}_{S, \sigma}^{\beta, l}( S^2) \bm{\rangle}$.

$ \textsf{ct}_{\alpha + \beta} = \textsf{ct}_{\alpha} + \textsf{ct}_{\beta}$

$ $

于是，以下性质成立：

$\phi_S(\textsf{ct}_{\alpha+\beta})\equiv(\Delta M^{\langle1\rangle}+E^{\langle1\rangle}+qK_1)(\Delta M^{\langle2\rangle}+E^{\langle2\rangle}+qK_2)+E_{\rm rl}\pmod Q$。

$ $




\item \textbf{\underline{重缩放}}

将 $\textsf{ct}_{\alpha + \beta} = (A_{\alpha + \beta}, B_{\alpha + \beta}) \bmod Q$ 更新为 $\textsf{ct'}_{\alpha + \beta} = \left(\left\lceil\dfrac{A_{\alpha + \beta}}{\Delta}\right\rfloor, \left\lceil\dfrac{B_{\alpha + \beta}}{\Delta}\right\rfloor\right) \bmod q$。

$ $

这一明文重缩放过程也可以视为密文 $\textsf{ct}_{\alpha + \beta}$ 按 $Q \rightarrow q$ 进行模数切换。


$ $

\end{enumerate}

注意，密文与密文相乘后，明文尺度因子 $\Delta=\left\lfloor\dfrac{q}{t}\right\rfloor$、密文模数 $q$ 和私钥 $S$ 均保持不变。



\end{tcolorbox}


\para{\textsf{ModRaise} 的目的：}在缩放之前把张量乘积模 $q$ 约减，可能丢失缩放舍入所需的信息。此教学变体改为在 $Q=q\Delta$ 下计算，因为被约减的 $Q$ 倍数在除以 $\Delta$ 后成为 $q$ 倍数。它并不要求每个明文乘积都小于 $Q$；对于多项式卷积及较大明文，这种断言一般不成立。仍需保留提升项 $qK_i$，并计入 \autoref{subsubsec:bfv-mult-cipher-rescaling} 推导的完整误差。

$ $

\para{交换\textsf{重线性化}与\textsf{重缩放}的顺序：}两种顺序都可以构造，但其求值密钥、模数和舍入误差界不同，不能直接使用同一组密钥和误差证明交换步骤。标准 FV 通常先缩放二次整数乘积，再在目标模数下重线性化。



\subsubsection{应用于 TFHE 的密文与密文相乘}
\label{subsubsec:bfv-mult-cipher-tfhe}

本小节有意采用另一种 LWE 相位 $b+\langle\vec a,\vec s\rangle$。若用于 \autoref{subsec:tfhe-enc-dec} 中相位为 $b-\langle\vec a,\vec s\rangle$ 的 TFHE 约定，须先将公开掩码取负，输出时再取负还原。张量乘积构造是一种分层 LWE 乘法替代方案，并非 TFHE 自举所使用的外积。它需要对乘积 $s_i s_j$ 的求值密钥、适当参数及完整重缩放误差界，并不能消除任意 LWE 乘法的噪声。


在 \autoref{subsec:tfhe-mult-cipher} 中，外积的输入是一个 LWE 密文和一个 GSW 密文。从任意 LWE 密文获得后者，可能需要电路自举。下述张量构造改为使用另一种二次规模的求值密钥，并要求相应的重缩放预算。这给出一种替代运算，并不宣称其性能优于 TFHE 外积或可编程自举。与本书 TFHE 的编码一样，要求 $t\mid q$ 且 $\Delta=q/t$。给定两个 LWE 密文，其整数提升满足：

$\textsf{LWE}_{\vec{s}, \sigma}(\Delta m^{\langle 1 \rangle} + e^{\langle 1 \rangle}) = (\vec{a}^{\langle 1 \rangle}, b^{\langle 1 \rangle }) \Rightarrow \,\,\, \vec{a}^{\langle 1 \rangle} \cdot \vec{s} + b^{\langle 1 \rangle} = \Delta m^{\langle 1 \rangle} + e^{\langle 1 \rangle} + k_1q$

$\textsf{LWE}_{\vec{s}, \sigma}(\Delta m^{\langle 2 \rangle} + e^{\langle 2 \rangle})  = (\vec{a}^{\langle 2 \rangle}, b^{\langle 2 \rangle}) \Rightarrow \,\,\, \vec{a}^{\langle 2 \rangle} \cdot \vec{s} + b^{\langle 2 \rangle} = \Delta m^{\langle 2 \rangle} + e^{\langle 2 \rangle} + k_2q$

$ $

将这两个等式相乘，得到以下关系：


$(\Delta m^{\langle 1 \rangle} + e^{\langle 1 \rangle} + k_1q) \cdot (\Delta m^{\langle 2 \rangle} + e^{\langle 2 \rangle} + k_2q)$

$ = (\vec{a}^{\langle 1 \rangle} \cdot \vec{s} + b^{\langle 1 \rangle}) \cdot (\vec{a}^{\langle 2 \rangle} \cdot \vec{s} + b^{\langle 2 \rangle})$ 

$ = b^{\langle 1 \rangle}b^{\langle 2 \rangle}  + (b^{\langle 2 \rangle}\vec{a}^{\langle 1 \rangle} + b^{\langle 1 \rangle}\vec{a}^{\langle 2 \rangle}) \cdot \vec{s} + (\vec{a}^{\langle 1 \rangle} \cdot \vec{s}) \cdot (\vec{a}^{\langle 2 \rangle} \cdot \vec{s}) $



$ = \underbrace{b^{\langle 1 \rangle}b^{\langle 2 \rangle}}_{d_0}  + \underbrace{(b^{\langle 2 \rangle}\vec{a}^{\langle 1 \rangle} + b^{\langle 1 \rangle}\vec{a}^{\langle 2 \rangle})}_{d_1} \cdot \vec{s} + \underbrace{(\vec{a}^{\langle 1 \rangle} \otimes \vec{a}^{\langle 2 \rangle})}_{\vec{d}_2} \cdot (\vec{s} \otimes \vec{s}) $  


$= d_0 + d_1\cdot \vec{s} + \vec{d}_2\cdot \vec{s}^{\otimes}$ \textcolor{red}{ $\rhd$ 其中 $\vec{s}^{\otimes} = \vec{s} \otimes \vec{s}$}

$ $

，其中，$\otimes$ 表示两个向量的外积。例如，给定两个长度为 $n$ 的向量 $\vec{v}$ 和 $\vec{u}$，$\vec{v} \otimes \vec{u}$ 等价于将以下 $n$ 个不同的、长度为 $n$ 的向量连接而成的长度为 $n^2$ 的向量：$v_0\cdot\vec{u}, \,\, v_1\cdot\vec{u}, \,\, \cdots v_{n-1}\cdot\vec{u}$。注意，上述关系与 BFV 密文与密文相乘中推导出的关系类似。因此，类似于 BFV 的其余步骤，可以对这一 LWE 项按如下方式进行重线性化与重缩放：


\begin{tcolorbox}[title={\textbf{\tboxlabel{\ref*{subsubsec:bfv-mult-cipher-tfhe}} 分层 LWE 张量乘法}}]


假设有以下两个 LWE 密文：

$\textsf{LWE}_{\vec{s}, \sigma}(\Delta M^{\langle 1 \rangle}) = (\vec{a}^{\langle 1 \rangle}, b^{\langle 1 \rangle}) \bmod q$，\text{ } 其中 $b^{\langle 1 \rangle} = -\vec{a}^{\langle 1 \rangle} \cdot \vec{s} + \Delta m^{\langle 1 \rangle} + e^{\langle 1 \rangle} \bmod q$

$\textsf{LWE}_{\vec{s}, \sigma}(\Delta M^{\langle 2 \rangle}) = (\vec{a}^{\langle 2 \rangle}, b^{\langle 2 \rangle}) \bmod q$，\text{ } 其中 $b^{\langle 2 \rangle} = -\vec{a}^{\langle 2 \rangle} \cdot \vec{s} + \Delta m^{\langle 2 \rangle} + e^{\langle 2 \rangle} \bmod q$

$ $

这两个 LWE 密文的乘法按如下步骤执行：

$ $

\begin{enumerate}

\item \textbf{\underline{模数提升（\textsf{ModRaise}）}}

将 $(\vec{a}^{\langle 1 \rangle}, b^{\langle 1 \rangle})$ 和 $(\vec{a}^{\langle 2 \rangle}, b^{\langle 2 \rangle})$ 各分量选定的整数提升视为模 $Q=q\Delta$ 的元素。其相位仍包含 $qk_1$ 和 $qk_2$，并不是在较大模数下对相同缩放明文的新鲜加密：

$(\vec{a}^{\langle 1 \rangle}, b^{\langle 1 \rangle}) \bmod Q$

$(\vec{a}^{\langle 2 \rangle}, b^{\langle 2 \rangle}) \bmod Q$


$ $

\item \textbf{\underline{乘法}}

在模 $Q$ 下计算以下多项式乘法：

$d_0 = b^{\langle 1 \rangle}b^{\langle 2 \rangle} \bmod Q$

$d_1 = b^{\langle 2 \rangle}\vec{a}^{\langle 1 \rangle} +  b^{\langle 1 \rangle}\vec{a}^{\langle 2 \rangle} \bmod Q$

$\vec{d}_2 = \vec{a}^{\langle 1 \rangle} \otimes \vec{a}^{\langle 2 \rangle} \bmod Q$


$ $

\item \textbf{\underline{重线性化}}


计算：

$ \textsf{ct}_\alpha = (d
_1, d_0)$

$ \textsf{ct}_\beta = \bm{\langle} \textsf{Decomp}^{\beta, l}(\vec{d}_2), \text{ } \textsf{Lev}_{\vec{s}, \sigma}^{\beta, l}(\vec{s}^{\otimes}) \bm{\rangle}$.

$ \textsf{ct}_{\alpha + \beta} = \textsf{ct}_{\alpha} + \textsf{ct}_{\beta}$

$ $

于是，以下性质成立：

$\phi_{+}(\textsf{ct}_{\alpha+\beta})\equiv(\Delta m^{\langle1\rangle}+e^{\langle1\rangle}+qk_1)(\Delta m^{\langle2\rangle}+e^{\langle2\rangle}+qk_2)+e_{\rm rl}\pmod Q$，其中 $\phi_+(\vec a,b)=b+\langle\vec a,\vec s\rangle$。

$ $




\item \textbf{\underline{重缩放}}

将 $\textsf{ct}_{\alpha + \beta} = (\vec{a}_{\alpha + \beta}, b_{\alpha + \beta}) \bmod Q$ 更新为 $\textsf{ct'}_{\alpha + \beta} = \left(\left\lceil\dfrac{\vec{a}_{\alpha + \beta}}{\Delta}\right\rfloor, \left\lceil\dfrac{b_{\alpha + \beta}}{\Delta}\right\rfloor\right) \bmod q$。

$ $

这一明文重缩放过程也可以视为密文 $\textsf{ct}_{\alpha + \beta}$ 按 $Q \rightarrow q$ 进行模数切换。


$ $

\end{enumerate}

注意，密文与密文相乘后，明文尺度因子 $\Delta=\dfrac{q}{t}$、密文模数 $q$ 和私钥 $\vec{s}$ 均保持不变。



\end{tcolorbox}



\textbf{求值密钥与条件性正确性。}设 LWE 维数为 $k$，选择整数分解权重 $g_j=\beta^j$，其中 $0\leq j<L$，并要求 $\beta^L\geq Q$。对每个有序下标对 $(u,v)$，公开相位为 $g_js_us_v+e_{u,v,j}^{\rm key}\pmod Q$ 的行密文。将 $\vec d_2$ 的每个分量在模 $Q$ 下精确分解为绝对值不超过 $D_g$ 的数字。若 $|e_{u,v,j}^{\rm key}|\leq B_{\rm key}$，则所得重线性化误差满足
\[
|e_{\rm rl}|\leq k^2LD_gB_{\rm key}.
\]
该密钥不同于 TFHE 的 GGSW 自举密钥。输入分量与明文均取中心代表，记 $h=\|\vec s\|_1$，并设 $|e_i|\leq B_i$，则整数相位中的进位满足
\[
|k_i|\leq\kappa_i:=\frac{h+2}{2}+\frac{B_i}{q}.
\]
若 $R_a,R_b$ 是最后除以 $\Delta$ 时各分量的舍入误差，直接展开可得
\[
\begin{aligned}
\phi_+(\textsf{ct}_{\rm out})&\equiv\Delta m_1m_2+\eta\pmod q,\\
\eta&=m_1e_2+m_2e_1+t(e_1k_2+e_2k_1)
      +\frac{e_1e_2+e_{\rm rl}}\Delta
      +\langle R_a,\vec s\rangle+R_b.
\end{aligned}
\]
其中 $qk_1m_2$、$qk_2m_1$ 和 $qt k_1k_2$ 在模 $q$ 下消失。特别地，由于 $t\mid q$，一般 BFV 教学变体中危险的模数余数项在此消失。解码的一个充分条件是
\[
\frac t2(B_1+B_2)+t(B_1\kappa_2+B_2\kappa_1)
+\frac{B_1B_2+k^2LD_gB_{\rm key}}\Delta+\frac{h+1}{2}
<\frac\Delta2.
\]
明文模 $t$ 约减只增加 $q$ 的倍数，不再增加误差。这给出了上述运算顺序的条件性证明。若改为先用 $1/\Delta=t/q$ 缩放整数张量，再在模 $q$ 下重线性化，则密钥模数为 $q$，求值密钥误差不再除以 $\Delta$，张量分量舍入的界为 $(1+h+h^2)/2$。这两种顺序各有独立论证，密钥不能互换。附带的精确检查使用非零输入误差和求值密钥误差验证了两种顺序，但不据此断言参数达到安全要求或性能更快。

\subsection{同态密钥切换}
\label{subsec:bfv-key-switching}

分解基及其加密各行必须使用相同约定。对于一般模数 $q$，可使用整数权重 $\beta^{i-1}$，并采用足够多的位数精确分解；若采用截断分解，残差与旧秘密密钥的乘积也必须计入密钥切换误差。

BFV 的密钥切换方案将 RLWE 密文的秘密密钥从 $S$ 改为 $S'$。其本质上是 RLWE 的密钥切换方案，只是在加密和解密公式中将 $A\cdot S$ 项的符号取反。具体而言，它等价于 $k = 1$ 时的另一种 GLWE 版本（\autoref{subsec:glwe-alternative}）的密钥切换方案（\autoref{sec:glwe-key-switching}），如下所示：

\begin{tcolorbox}[title={\textbf{\tboxlabel{\ref*{subsec:bfv-key-switching}} BFV 的密钥切换}}]
$\textsf{RLWE}_{S',\sigma}(\Delta M) = (0, B) + \bm{\langle} \textsf{Decomp}^{\beta, l}(A), \text{ } \textsf{RLev}_{S', \sigma}^{\beta, l}(S) \bm{\rangle}$
\end{tcolorbox}




\subsection{输入向量槽位的同态旋转}
\label{subsec:bfv-rotation}


本节将介绍：当输入向量 $\vec{v}$ 已被编码为多项式并加密为 RLWE 密文后，如何对其元素进行同态旋转。在 \autoref{subsec:coeff-rotation} 中，我们学习了如何旋转多项式的系数。然而，旋转明文多项式 $M(X)$ 或 RLWE 密文多项式 $(A(X), B(X))$，不一定会旋转生成它们的输入向量。

对输入向量槽位（即输入向量）进行同态旋转的关键要求是：该运算必须作用于 RLWE 密文，并且在运算结束后，对 RLWE 密文进行解密和解码所恢复出的输入向量，应处于预期的旋转状态。我们将这一任务分为以下两个子问题：

\begin{enumerate}[leftmargin=3\parindent]
\item 如何通过将明文多项式 $M$ 更新为 $M'$，间接旋转输入向量？

\item 如何通过将 RLWE 密文多项式 $(A, B)$ 更新为 $(A', B')$，间接将明文多项式 $M$ 更新为 $M'$？
\end{enumerate}



\subsubsection{通过更新明文多项式旋转输入向量槽位}

在这一任务中，目标是修改明文多项式 $M(X)$，使输入向量 $\vec{v}$ 的前半部分元素在其内部向左循环移动 $h$ 个位置，同时使 $\vec{v}$ 的后半部分元素也在其内部向左循环移动 $h$ 个位置。具体而言，如果 $\vec{v}$ 定义如下：

$\vec{v} = (v_0, v_1, \cdots, v_{n-1})$

$ $

那么，将移动 $h$ 个位置后的向量 $\vec{v}^{\langle h \rangle}$ 记为：

$\vec{v}^{\langle h \rangle} = (\underbrace{v_h, v_{h+1}, \cdots, v_0, v_1, \cdots, v_{h-2}, v_{h-1},}_{\text{前半部分 $\dfrac{n}{2}$ 个元素向左旋转 $h$ 个位置}} \text{ } \underbrace{v_{\frac{n}{2}+h}, v_{\frac{n}{2}+h+1}, \cdots, v_{\frac{n}{2}+h-2}, v_{\frac{n}{2}+h-1}}_{\text{后半部分 $\dfrac{n}{2}$ 个元素向左旋转 $h$ 个位置}})$

$ $

回顾 \autoref{subsec:bfv-batch-encoding}，BFV 编码方案的组成部分如下：

$\vec{v} = (v_0, v_1, v_2, \cdots, v_{n-1})$ \textcolor{red}{ $\rhd$ 包含 $n$ 个槽位的输入向量}

$ $

$W =  \begin{bmatrix}
1 & 1 & 1 & \cdots & 1\\
(\omega) & (\omega^3) & (\omega^5) & \cdots & (\omega^{2n-1})\\
(\omega)^2 & (\omega^3)^2 & (\omega^5)^2 & \cdots & (\omega^{2n-1})^2\\
\vdots & \vdots & \vdots & \ddots & \vdots \\
(\omega)^{n-1} & (\omega^3)^{n-1} & (\omega^5)^{n-1} & \cdots & (\omega^{2n-1})^{n-1}\\
\end{bmatrix}$

$= \begin{bmatrix}
1 & 1 & \cdots & 1 & 1 & \cdots & 1 & 1\\
(\omega) & (\omega^3) & \cdots & (\omega^{n - 1}) & (\omega^{-(n - 1)}) & \cdots & (\omega^{-3}) & (\omega^{-1})\\
(\omega)^2 & (\omega^3)^2 & \cdots & (\omega^{n - 1})^2 & (\omega^{-(n - 1)})^2 & \cdots & (\omega^{-3})^2 & (\omega^{-1})^2\\
\vdots & \vdots & \vdots & \ddots & \vdots \\
(\omega)^{n-1} & (\omega^3)^{n-1} & \cdots & (\omega^{n - 1})^{n-1} & (\omega^{-(n - 1)})^{n-1} & \cdots & (\omega^{-3})^{n-1} & (\omega^{-1})^{n-1}\\
\end{bmatrix}$

，其中 $\omega = g^{\frac{t - 1}{2n}} \bmod t$（$g$ 是 $\mathbb{Z}_t^{\times}$ 的一个生成元）

\textcolor{red}{ $\rhd$ 将 $\vec{v}$ 转换为 $\vec{m}$ 的编码矩阵（即明文多项式 $M(X)$ 的 $n$ 个系数）}
 


$ $

$\Delta\vec{m}  = n^{-1}\cdot\Delta W\cdot I_n^R\cdot\vec{v}$  

\textcolor{red}{ $\rhd$ 包含明文多项式经缩放后的 $n$ 个整数系数的向量}

$ $

$\Delta M = \sum\limits_{i=0}^{n-1} ( \Delta m_i  X^i)$ 

\textcolor{red}{ $\rhd$ 对输入向量 $\vec{v}$ 进行同构编码得到的整数多项式}


$ $

由 \autoref{subsec:poly-vector-transformation} 可知，对多项式 $M(X)$ 进行解码 $\vec{v}$，等价于在以下 $n$ 个不同的本原 $(\mu=2n)$ 次单位根处计算 $M(X)$ 的值：$\{\omega, \omega^3, \omega^5, \cdots, \omega^{2n-3}, \omega^{2n-1} \}$。因此，上述解码过程等价于：

$\vec{v} = \dfrac{W^T \Delta \vec{m}}{\Delta} = \bm{(}W^T_0\cdot \vec{m}, \text{ } W^T_1\cdot \vec{m}, \text{ } W^T_2\cdot \vec{m}, \text{ } \cdots, \text{ } W^T_{n-1}\cdot \vec{m}\bm{)}$  

$= \bm{(} \text{ } M(\omega), \text{ } M(\omega^3), \text{ } M(\omega^5), \cdots, M(\omega^{2n-3}), \text{ } M(\omega^{2n-1}) \text{ } \bm{)}$


$= \bm{(} \text{ } M(\omega), \text{ } M(\omega^3), \text{ } M(\omega^5), \cdots, M(\omega^{n-3}), \text{ } M(\omega^{n-1}), \text{ } M(\omega^{-(n-1)}), \text{ } M(\omega^{-(n-3)}), \cdots, M(\omega^{-3}), \text{ } M(\omega^{-1}) \text{ } \bm{)}$


%%$\textsf{LWE}_{S, \sigma}(\Delta M) = (A, B) \in \mathcal{R}_{\langle n, q \rangle }^{2}$ \textcolor{red}{  $\rhd$ An LWE ciphertext that encrypts $\Delta M$}

%$A\cdot S + B = \Delta M + E$


$ $

接下来的任务是将 $M(X)$ 修改为 $M'(X)$，使解码 $M'(X)$ 得到经过修改的输入向量 $\vec{v}^{\langle h \rangle}$：它将 $\vec{v}$ 的前半部分元素在其内部向左循环旋转 $h$ 个位置，并将后半部分元素也在其内部向左循环旋转 $h$ 个位置。为实现这一旋转，我们采用以下两步方案：

\begin{enumerate}[leftmargin=3\parindent]
\item 为将 $M(X)$ 转换为 $M'(X)$，定义新的映射 $\sigma_M$ 如下：

$\sigma_M : (M(X), h) \in (\mathcal{R}_{\langle n, t \rangle}, \mathbb{Z}_{n}) \longrightarrow M'(X) \in \mathcal{R}_{\langle n, t \rangle}$

，其中，$h$ 是要对 $\vec{v}$ 施加的旋转位置数。
\item 为将 $M'(X)$ 解码为旋转后的输入向量 $\vec{v}^{\langle h \rangle}$，需要修改 \textsf{Encoding\textsubscript{1}}（\autoref{subsubsec:bfv-encoding-1}）中的同构映射 $\sigma : M(X) \in \mathcal{R}_{\langle n, t \rangle} \longrightarrow \vec{v} \in \mathbb{Z}^n$，重新设计解码方案。
\end{enumerate}

$ $

\para{将 $\bm{M(X)}$ 转换为 $\bm{M'(X)}$：} 第一个任务是将 $M(X)$ 转换为 $M'(X)$，等价于应用新的映射 $\sigma_M : (M(X), h) \in (\mathcal{R}_{\langle n, t \rangle}, \mathbb{Z}_{n}) \longrightarrow M'(X) \in \mathcal{R}_{\langle n, t \rangle}$，使解码 $M'(X)$ 得到旋转后的输入向量 $\vec{v}^{\langle h \rangle}$：$\vec{v}$ 的前半部分元素在其内部向左循环旋转 $h$ 个位置，后半部分元素也在其内部向左循环旋转 $h$ 个位置。为了设计满足这一要求的 $\sigma_M$，我们将使用具有以下两个特殊性质的数 $5^j$（这些性质来自数论）：
\begin{itemize}
\item 当整数 $j$ 满足 $0 \leq j < \dfrac{n}{2}$ 时，$(5^j \bmod 2n)$ 和 $(-5^j \bmod 2n)$ 生成 $[0, 2n)$ 之间的所有奇数。
\item 对每个满足 $0 \leq j < \dfrac{n}{2}$ 的整数 $j$，均有 $(5^j \bmod 2n) + (-5^j \bmod 2n) \equiv 0 \bmod 2n$。
\end{itemize}
例如，假设模数 $2n = 16$。那么，

\begin{multicols}{2}
\noindent $5^0 \bmod 16 = 1$
\newline $5^1 \bmod 16 = 5$
\newline $5^2 \bmod 16 = 9$
\newline $5^3 \bmod 16 = 13$
\newline $-(5)^0 \bmod 16 = 15$
\newline $-(5)^1 \bmod 16 = 11$
\newline $-(5)^2 \bmod 16 = 7$
\newline $-(5)^3 \bmod 16 = 3$
\end{multicols}


如上所示，$0 \leq j < 4$ 生成 $[0, 16)$ 之间的所有奇数。此外，对 $0 \leq j < 4$ 中的每个 $j$，均有 $(5^j \bmod 16) + (-5^j \bmod 16) = 16$。

$ $

定义 $J(h) = 5^h \bmod 2n$，\text{ } 以及 \text{ } $J_*(h) = -5^h \bmod 2n$。基于 $J(h)$ 和 $J_*(h)$，将映射 $\sigma_M : M(X) \rightarrow M'(X)$ 定义如下：

$\sigma_{M} : M(X) \rightarrow M(X^{J(h)})$

$ $

给定明文多项式 $M(X)$，为了使其解码得到的输入向量 $\vec{v}$ 的前半部分元素向左循环旋转 $h$ 个位置，后半部分元素也向左循环旋转 $h$ 个位置，我们通过应用 $\sigma_M$ 映射，将当前明文多项式 $M(X)$ 更新为新的多项式 $M'(X) = M(X^{J(h)}) = M(X^{5^h})$，其中，$h$ 是 $\vec{v}$ 的前、后半部分元素各自向左旋转的位置数。


$ $

\para{将 $\bm{M'(X)}$ 解码为 \boldmath$\vec{v}^{\langle h \rangle}$}：第二个任务是修改原始解码方案，使其能够成功地将 $M'(X)$ 解码为旋转后的输入向量 $\vec{v}^{\langle h \rangle}$。为此，我们将原来的同构映射 $\sigma : M(X) \longrightarrow \vec{v}$ 从：

$\sigma: M(X) \in \mathcal{R}_{\langle n, q \rangle} \longrightarrow  \bm{(}M(\omega), M(\omega^3),M(\omega^5), \cdots, M(\omega^{2n-1})\bm{)} \in \mathbb{Z}^n$ \textcolor{red}{ \text{ } \# 该映射在 \autoref{subsec:poly-vector-transformation} 中设计}

$ $

，改为：

$\sigma_J: M(X) \in \mathcal{R}_{\langle n, t \rangle} \longrightarrow  \bm{(}M(\omega^{J(0)}),M(\omega^{J(1)}),M(\omega^{J(2)}), \cdots, M(\omega^{J(\frac{n}{2} - 1)}), $

\textcolor{white}{$\sigma_J: M(X) \in \mathcal{R}_{\langle n, q \rangle} \longrightarrow  \bm{(}$} $M(\omega^{J_*(0)}), \text{ } M(\omega^{J_*(1)}), \text{ } M(\omega^{J_*(2)}), \cdots, M(\omega^{J_*(\frac{n}{2} - 1)}) \text{ } \bm{)} \in \mathbb{Z}^{n}$


$ $

$\sigma$ 与 $\sigma_J$ 的共同点在于：它们都在 $n$ 个不同的本原 $(\mu=2n)$ 次单位根处计算多项式 $M(X)$ 的值（即 $w^i$，其中 $i$ 取 $[0,2n]$ 之间的所有奇数）。对于 $\sigma_J$ 映射，注意当 $j$ 遍历 $0 \leq j < \dfrac{n}{2} $ \text{ } 时，$J(j) = 5^j \bmod 2n$ \text{ } 和 \text{ } $J_*(j) = -5^j \bmod 2n$ \text{ } 覆盖了 $[0, 2n]$ 之间的所有奇数。因此，在 $0 \leq j < \frac{n}{2}$ 范围内，$\omega^{J(j)}$ 和 $\omega^{J_*(j)}$ 覆盖了全部 $n$ 个不同的本原 $(\mu=2n)$ 次单位根。

同时，$\sigma$ 和 $\sigma_J$ 的差别在于输出向量元素的顺序。在 $\sigma$ 映射中，对 $M(X)$ 求值所用坐标的顺序为 $\omega, \omega^3, \cdots, \omega^{2n - 1}$；而在 $\sigma_J$ 映射中，求值所用坐标的顺序为 $\omega^{J(0)}, \omega^{J(1)}, \cdots, \omega^{J(\frac{n}{2} - 1)}, \omega^{J_*(0)}, \omega^{J_*(1)}, \cdots, \omega^{J_*(\frac{n}{2} - 1)}$。稍后将说明为何这样修改顺序。

在原始的 \textsf{Decoding\textsubscript{2}} 过程（\autoref{subsec:bfv-batch-encoding}）中，对明文多项式 $M(X)$ 应用 $\sigma$ 映射，等价于进行以下计算：

$\vec{v}_{'} = \bm{(} \text{ } M(\omega), \text{ } M(\omega^3), \text{ } M(\omega^5), \cdots, M(\omega^{2n-3}), \text{ } M(\omega^{2n-1}) \bm{)}$ 

$ $

$ = \bm{(} \text{ } M(\omega), \text{ } M(\omega^3), \text{ } M(\omega^5), \cdots, M(\omega^{n-1}), \text{ } M(\omega^{-(n-1)}), \cdots, \text{ } M(\omega^{-3}), \text{ } M(\omega^{-1}) \text{ } \bm{)}$ 

$ $

$= \bm{(} \text{ } W^T_0\cdot \vec{m}, \text{ } W^T_1\cdot \vec{m}, \text{ } W^T_2\cdot \vec{m}, \text{ } \cdots, \text{ } W^T_{n-1}\cdot \vec{m} \text{ } \bm{)} $ 
$ $ \textcolor{red}{ $\rhd$ 其中，$W_i^T$ 是 $W^T$ 的第 $(i+1)$ 行}

$ = W^T \vec{m} $
%$= \dfrac{I_n^R}{n}\cdot \bm{(} \text{ } M(\omega), \text{ } M(\omega^3), \text{ } M(\omega^5), \cdots, M(\omega^{2n-3}), \text{ } M(\omega^{2n-1}) \text{ } )$




$ $

类似地，对明文多项式 $M(X)$ 应用修改后的 $\sigma_J$ 映射，等价于进行以下计算：

$\vec{v}_{'} =  \bm{(} \text{ } 
M(\omega^{J(0)}), \text{ } M(\omega^{J(1)}), \text{ } M(\omega^{J(2)}), \cdots,  M(\omega^{J(\frac{n}{2}-1)}), \text{ } M(\omega^{J_*(0)}), \text{ } M(\omega^{J_*(1)}), \cdots,  M(\omega^{J_*(\frac{n}{2}-1)}) \text{ } \bm{)}$


$ $

$= \bm{(} \text{ } \hathat{W}^*_0\cdot \vec{m}, \text{ } \hathat{W}^*_1\cdot \vec{m}, \text{ } \hathat{W}^*_2\cdot \vec{m}, \text{ } \cdots, \text{ } \hathat{W}^*_{n-1}\cdot \vec{m} \text{ } \bm{)} $ 

$ $

$ = \hathat{W}^* \vec{m} $，其中

$ $

\noindent $\hathat{W}^* = \begin{bmatrix}
1 & (\omega^{J(0)}) & (\omega^{J(0)})^2 & \cdots & (\omega^{J(0)})^{n-1}\\
1 & (\omega^{J(1)}) & (\omega^{J(1)})^2 & \cdots & (\omega^{J(1)})^{n-1}\\
1 & (\omega^{J(2)}) & (\omega^{J(2)})^2 & \cdots & (\omega^{J(2)})^{n-1}\\
\vdots & \vdots & \vdots & \ddots & \vdots \\
1 & (\omega^{J(\frac{n}{2}-1)}) & (\omega^{J(\frac{n}{2}-1)})^2 & \cdots & (\omega^{J(\frac{n}{2}-1)})^{n-1}\\
1 & (\omega^{J_*(0)}) & (\omega^{J_*(0)})^2 & \cdots & (\omega^{J_*(0)})^{n-1}\\
1 & (\omega^{J_*(1)}) & (\omega^{J_*(1)})^2 & \cdots & (\omega^{J_*(1)})^{n-1}\\
1 & (\omega^{J_*(2)}) & (\omega^{J_*(2)})^2 & \cdots & (\omega^{J_*(2)})^{n-1}\\
\vdots & \vdots & \vdots & \ddots & \vdots \\
1 & (\omega^{J_*(\frac{n}{2}-1)}) & (\omega^{J_*(\frac{n}{2}-1)})^2 & \cdots & (\omega^{J_*(\frac{n}{2}-1)})^{n-1}\\
\end{bmatrix}$

$ $

$ $

\noindent $\hathat{W} = \begin{bmatrix}
1 & 1 & \cdots & 1 & 1 & 1 & \cdots & 1\\
(\omega^{J(\frac{n}{2} - 1)}) & (\omega^{J(\frac{n}{2} - 2)}) & \cdots & (\omega^{J(0)}) & (\omega^{J_*(\frac{n}{2} - 1)}) & (\omega^{J_*(\frac{n}{2} - 2)}) & \cdots & (\omega^{J_*(0)})\\
(\omega^{J(\frac{n}{2} - 1)})^2 & (\omega^{J(\frac{n}{2} - 2)})^2 & \cdots & (\omega^{J(0)})^2 & (\omega^{J_*(\frac{n}{2} - 1)})^2 & (\omega^{J_*(\frac{n}{2} - 2)})^2 & \cdots & (\omega^{J_*(0)})^2 \\
\vdots & \vdots & \ddots & \vdots & \vdots & \vdots & \vdots & \vdots \\
(\omega^{J(\frac{n}{2} - 1)})^{n-1} & (\omega^{J(\frac{n}{2} - 2)})^{n-1} & \cdots & (\omega^{J(0)})^{n-1} & (\omega^{J_*(\frac{n}{2} - 1)})^{n-1} & (\omega^{J_*(\frac{n}{2} - 2)})^{n-1} & \cdots  & (\omega^{J_*(0)})^{n-1}
\end{bmatrix}$

$ $

$ $

从现在开始，将 \textsf{Encoding\textsubscript{1}} 过程（\autoref{subsec:bfv-batch-encoding}）中的 $W$ 替换为 $\hathat{W}$，并将 \textsf{Decoding\textsubscript{2}} 过程中的 $W^T$ 替换为 $\hathat{W}^*$。

$ $

为证明 $\hathat{W}$ 和 $W$ 一样是有效的编码矩阵，且 $\hathat{W}^*$ 和 $W^T$ 一样是有效的解码矩阵，需要证明以下两点：
\begin{itemize}[leftmargin=3\parindent]

\item \textbf{${\hathat{\bm W}}$ 是 $\bm{n}$ 维向量空间的一组基：} 这一点成立，因为 $\hathat{W}$ 只是对 $W$ 的列重新排序，仍然是 $\bm{n}$ 维向量空间的一组基。

\item \boldmath{$\hathat{W}^{*} \cdot \hathat{W} = n \cdot I_n^R$（满足 \autoref{subsec:vandermonde-euler} 中的定理~\ref*{subsec:vandermonde-euler}）}：这一证明分为两个子证明：
\begin{enumerate}[leftmargin=2\parindent]
\item \boldmath{${\hathat{ W}^{ *} \cdot \hathat{ W}}$ 的反对角线上各值为 $\bm{n}$：} \boldmath$\hathat{ W}^{ *} \cdot \hathat{ W}$ 的反对角线上每个元素的计算结果均为 $\sum\limits_{i=0}^{n-1} \omega^{2nik} = n$，其中 $k$ 为某个整数。
\item \boldmath{${\hathat{ W}^{ *} \cdot \hathat{ W}}$ 的其余位置各值均为 0：} 除反对角线上的元素外，其余所有元素均为 $\sum\limits_{i=0}^{n-1}r^i=\dfrac{r^n-1}{r-1}=0\quad(r=\omega^{a+b}\ne1)$（由等比数列求和公式可得）。
\end{enumerate}

我们提供了 \href{https://github.com/wang2000star/fhe-textbook-cn/blob/chinese-edition/source%20code/bfv_j_matrix_inverse_proof.py}{\underline{Python 脚本}}，用于通过实验验证这一点。

\end{itemize}




因此，$\hathat{W}^*$ 和 $\hathat{W}$ 是将 $\vec{v}$ 转换为 $M(X)$ 的有效编码与解码矩阵。

$ $

现在考虑 $\vec{v}^{\langle h \rangle}$ 的结构（即将 $\vec{v}$ 的前半部分元素在其内部向左循环旋转 $h$ 个位置，同时将后半部分元素也在其内部向左循环旋转 $h$ 个位置）。回顾 $\vec{v}$ 的表达式：

$\vec{v} = \bm{(} \text{ } 
M(\omega^{J(0)}), \text{ } M(\omega^{J(1)}), \cdots,  M(\omega^{J(\frac{n}{2}-1)}), \text{ } 
M(\omega^{J_*(0)}), \text{ } M(\omega^{J_*(1)}),  \cdots,  M(\omega^{J_*(\frac{n}{2}-1)}) \bm{)}$


\text{ } \text{ } $= \bm{(} \text{ } \hathat{W}^*_0\cdot \vec{m}, \text{ } \hathat{W}^*_1\cdot \vec{m}, \text{ } \hathat{W}^*_2\cdot \vec{m}, \text{ } \cdots, \text{ } \hathat{W}^*_{n - 1}\cdot \vec{m} \text{ } \bm{)} $ 

$ $

因此，分别将 $\vec{v}$ 的前、后半部分元素向左旋转 $h$ 个位置，所得 $\vec{v}^{\langle h \rangle}$ 的状态如下：

$\vec{v}^{\langle h \rangle} = \bm{(} \text{ } M(\omega^{J(h)}), \text{ } M(\omega^{J(h+1)}), \cdots, M(\omega^{J(\frac{n}{2}-1)}), \text{ } M(\omega^{J(0)}), \text{ }M(\omega^{J(1)}), \cdots, \text{ } M(\omega^{J(h-2)}), \text{ } M(\omega^{J(h-1)}),$

\text{ }\text{ }\text{ } $ M(\omega^{J_*(h)}), \text{ } M(\omega^{J_*(h+1)}), \cdots, M(\omega^{J_*(\frac{n}{2}-1)}), \text{ } M(\omega^{J_*(0)}), \text{ }M(\omega^{J_*(1)}), \cdots, \text{ } M(\omega^{J_*(h-2)}), \text{ } M(\omega^{J_*(h-1)}) \text{ } \bm{)}$

$ $

注意，上述 $\vec{v}^{\langle h \rangle}$ 的计算等价于：将 $\hathat{W}^*$ 的上方 $\frac{n}{2}$ 行在其内部向上循环旋转 $h$ 个位置，将 $\hathat{W}^{*}$ 的下方 $\frac{n}{2}$ 行在其内部向上循环旋转 $h$ 个位置，再将得到的矩阵与 $\vec{m}$ 相乘。然而，实际中不希望如此直接修改解码矩阵 $\hathat{W}^*$，因为这样解码矩阵就无法保持一致。因此，我们将通过把 $\vec{m}$ 修改为 $\vec{m}_{'}$（即把 $M(X)$ 修改为 $M'(X)$）来保持关系 $\vec{v}^{\langle h \rangle} = \hathat{W}^* \cdot \vec{m}_{'}$ 成立，而不直接修改 $\hathat{W}^*$。提取 $\hathat{W}^*$ 的上半部分行，将这一 $\dfrac{n}{2} \times n$ 矩阵记为 $\hathat{H}_1^*$。那么，$\hathat{H}_1^*$ 等价于形式为 $V(\omega^{J(0)}, \omega^{J(1)}, \cdots, \omega^{J(\frac{n}{2}-1)})$ 的范德蒙德矩阵（\autoref{subsec:vandermonde} 中的定义~\ref*{subsec:vandermonde}）。类似地，提取 $\hathat{W}^*$ 的下半部分行，将这一 $\dfrac{n}{2} \times n$ 矩阵记为 $\hathat{H}_2^*$。那么，$\hathat{H}_2^*$ 等价于形式为 $V(\omega^{J_*(0)}, \omega^{J_*(1)}, \cdots, \omega^{J_*(\frac{n}{2}-1)})$ 的范德蒙德矩阵（\autoref{subsec:vandermonde} 中的定义~\ref*{subsec:vandermonde}）。

现在，将 $\hathat{H}_1^*$ 的各行向上旋转 $h$ 个位置，并记为 $\hathat{H}_1^{*\langle h \rangle}$；将 $\hathat{H}_2^*$ 的各行向上旋转 $h$ 个位置，并记为 $\hathat{H}_2^{*\langle h \rangle}$。将由 $\vec{v}^{\langle h \rangle}$ 的前半部分元素构成的 $\dfrac{n}{2}$ 维向量记为 $\vec{v}_1^{\langle h \rangle}$，将由 $\vec{v}^{\langle h \rangle}$ 的后半部分元素构成的 $\dfrac{n}{2}$ 维向量记为 $\vec{v}_2^{\langle h \rangle}$。
那么，计算（即解码）$\vec{v}_1^{\langle h \rangle} = \hathat{H}_1^{*\langle h \rangle}\cdot \vec{m}$，等价于将 $M(X)$ 修改为 $M'(X) = M(X^{J(h)})$（其系数向量为 $\vec{m}^{\langle h \rangle}$）后，再计算（即解码）$\vec{v}_1^{\langle h \rangle} = \hathat{H}_1^{*}\cdot \vec{m}^{\langle h \rangle}$。原因如下：

$\vec{v}_1^{\langle h \rangle} =\hathat{H}_1^{*\langle h \rangle}\cdot \vec{m}$

$ = \bm{(} \text{ } \hathat{W}^*_{h}\cdot\vec{m}, \text{ } \hathat{W}^*_{h+1}\cdot\vec{m}, \text{ } \hathat{W}^*_{h+2}\cdot\vec{m}, \cdots, \hathat{W}^*_{\frac{n}{2}-1}\cdot\vec{m}, \text{ } \hathat{W}^*_{0}\cdot\vec{m}, \text{ }\hathat{W}^*_{1}\cdot\vec{m}, \cdots, \text{ } \hathat{W}^*_{h-2}\cdot\vec{m}, \text{ } \hathat{W}^*_{h-1}\cdot\vec{m}  \bm{)}$

$ $

$= \bm{(} \text{ } M((\omega^{J(h)})^{J(0)}), \text{ } M((\omega^{J(h)})^{J(1)}), \text{ } M((\omega^{J(h)})^{J(2)}), \cdots, M((\omega^{J(h)})^{J(\frac{n}{2}-1)}) \text{ } \bm{)}$


\textcolor{red}{ $\rhd$ 这等价于在以下 $\dfrac{n}{2}$ 个不同的 $(\mu=2n)$ 次单位根处计算多项式 $M(X^{J(h)})$ 的值：$\omega^{J(0)}, \omega^{J(1)},  \omega^{J(2)}, \cdots, \omega^{J(\frac{n}{2} - 1)}$}

$ $

$= \bm{(} \text{ } M(\omega^{J(h)\cdot J(0)}), \text{ } M(\omega^{J(h)\cdot J(1)}), \text{ } M(\omega^{J(h)\cdot J(2)}), \cdots, M(\omega^{J(h)\cdot J(\frac{n}{2}-1)}) \text{ } \bm{)}$

$= \bm{(} \text{ } M(\omega^{5^h\cdot 5^0}), \text{ } M(\omega^{5^h\cdot 5^1}), \text{ } M(\omega^{5^h\cdot 5^2}), \cdots, M(\omega^{5^h\cdot 5^{n/2-1}}) \text{ } \bm{)}$

$= \bm{(} \text{ } M(\omega^{5^{h}}), \text{ } M(\omega^{5^{h+1}}), \text{ } M(\omega^{5^{h+2}}), \cdots, M(\omega^{5^{h+n/2-1}}) \text{ } \bm{)}$ \textcolor{red}{ $\rhd$ 注意 $5^{\frac{n}{2}} \bmod 2n = 1$}

$ $

$= \bm{(} \text{ } M(\omega^{J(h)}), \text{ } M(\omega^{J(h+1)}), \text{ } M(\omega^{J(h+2)}), \cdots, M(\omega^{J(\frac{n}{2}-1)}), \text{ } M(\omega^{J(0)}), \text{ }M(\omega^{J(1)}),$ 
                                    
$\cdots, \text{ } M(\omega^{J(h-2)}), \text{ } M(\omega^{J(h-1)})  \bm{)}$ 

$ $

$= \hathat{H}_1^{*}\cdot \vec{m}^{\langle h \rangle}$ \textcolor{red}{ \text{ } \# 其中，$\vec{m}^{\langle h \rangle}$ 包含多项式 $M(X^{J(h)})$ 的 $n$ 个系数 }


$ $

类似地，计算（即解码）$\vec{v}_2^{\langle h \rangle} = \hathat{H}_2^{*\langle h \rangle}\cdot \vec{m}$，等价于将 $M(X)$ 修改为 $M'(X) = M(X^{J(h)})$（其系数向量为 $\vec{m}^{\langle h \rangle}$）后，再计算（即解码）$\vec{v}_2^{\langle h \rangle} = \hathat{H}_2^{*}\cdot \vec{m}^{\langle h \rangle}$。原因如下：

$\vec{v}_2^{\langle h \rangle} =\hathat{H}_2^{*\langle h \rangle}\cdot \vec{m}$

$ $


$ = \bm{(} \text{ } \hathat{W}^*_{\frac{n}{2} +h}\cdot\vec{m}, \text{ } \hathat{W}^*_{\frac{n}{2} + h+1}\cdot\vec{m}, \text{ } \hathat{W}^*_{\frac{n}{2} +h+2}\cdot\vec{m}, \cdots, \hathat{W}^*_{n -1}\cdot\vec{m}, \text{ } \hathat{W}^*_{\frac{n}{2}}\cdot\vec{m}, \text{ }\hathat{W}^*_{\frac{n}{2} + 1}\cdot\vec{m}, \cdots,$

\text{ }\text{ }\text{ }\text{ }$ \text{ } \hathat{W}^*_{\frac{n}{2} + h-2}\cdot\vec{m}, \text{ } \hathat{W}^*_{\frac{n}{2} + h-1}\cdot\vec{m}  \bm{)}$

$ $

$= \bm{(} \text{ } M((\omega^{J(h)})^{J_*(0)}), \text{ } M((\omega^{J(h)})^{J_*(1)}), \text{ } M((\omega^{J(h)})^{J_*(2)}), \cdots, M((\omega^{J(h)})^{J_*(\frac{n}{2}-1)}) \text{ } \bm{)}$


\textcolor{red}{ $\rhd$ 这等价于在以下 $\dfrac{n}{2}$ 个不同的 $(\mu=2n)$ 次单位根处计算多项式 $M(X^{J(h)})$ 的值：$\omega^{J_*(0)}, \omega^{J_*(1)},  \omega^{J_*(2)}, \cdots, \omega^{J_*(\frac{n}{2} - 1)}$}

$ $

$= \bm{(} \text{ } M(\omega^{J(h)\cdot J_*(0)}), \text{ } M(\omega^{J(h)\cdot J_*(1)}), \text{ } M(\omega^{J(h)\cdot J_*(2)}), \cdots, M(\omega^{J(h)\cdot J_*(\frac{n}{2}-1)}) \text{ } \bm{)}$

$= \bm{(} \text{ } M(\omega^{5^h\cdot -5^0}), \text{ } M(\omega^{5^h\cdot -5^1}), \text{ } M(\omega^{5^h\cdot -5^2}), \cdots, M(\omega^{5^h\cdot -5^{n/2-1}}) \text{ } \bm{)}$

$= \bm{(} \text{ } M(\omega^{-5^{h}}), \text{ } M(\omega^{-5^{h+1}}), \text{ } M(\omega^{-5^{h+2}}), \cdots, M(\omega^{-5^{h+n/2-1}}) \text{ } \bm{)}$ \textcolor{red}{ $\rhd$ 注意 $-(5^{\frac{n}{2}} \bmod 2n) = -1$}

$ $

$= \bm{(} \text{ } M(\omega^{J_*(h)}), \text{ } M(\omega^{J_*(h+1)}), \text{ } M(\omega^{J_*(h+2)}), \cdots, M(\omega^{J_*(\frac{n}{2}-1)}), \text{ } M(\omega^{J_*(0)}), \text{ }M(\omega^{J_*(1)}),$ 
                                    
$\cdots, \text{ } M(\omega^{J_*(h-2)}), \text{ } M(\omega^{J_*(h-1)})  \bm{)}$ 

$ $

$= \hathat{H}_2^{*}\cdot \vec{m}^{\langle h \rangle}$


$ $

上述推导表明，$\vec{v}_1^{\langle h \rangle} = \hathat{H}_1^{*}\cdot \vec{m}^{\langle h \rangle}$，且 $\vec{v}_2^{\langle h \rangle} = \hathat{H}_2^{*}\cdot \vec{m}^{\langle h \rangle}$。结合这两点，得到结论 $\vec{v}^{\langle h \rangle} = \hathat{W}^* \cdot \vec{m}^{\langle h \rangle}$：将输入向量 $\vec{v}$ 的前半部分元素向左旋转 $h$ 个位置、后半部分元素也向左旋转 $h$ 个位置，等价于将明文多项式 $M(X)$ 更新为 $M(X^{J(h)})$，再使用解码矩阵 $\hathat{W}^*$ 进行解码。



然而，现在出现了一个新问题：无法直接将明文 $M(X)$ 更新为 $M(X^{J(h)})$，因为 $M(X)$ 已被加密为 RLWE 密文。因此，需要更新 RLWE 密文分量 $(A, B)$，从而\textit{间接地}将 $M(X)$ 更新为 $M(X^{J(h)})$。下一小节将介绍具体方法。


\subsubsection{通过更新密文多项式来更新明文多项式}

给定 RLWE 密文 $\textsf{ct} = (A, B)$，目标是将 $\textsf{ct} = (A, B)$ 更新为 $C^{\langle h \rangle} = (A^{\langle h \rangle}, B^{\langle h \rangle})$，使其解密后得到输入向量 $\vec{v}^{\langle h \rangle}$。也就是说，应满足以下关系：

$\textsf{RLWE}^{-1}_{S, \sigma}\bm{(} \text{ } C^{\langle h \rangle} = (A^{\langle h \rangle}, B^{\langle h \rangle}) \text{ }\bm{)} = \Delta M(X^{J(h)}) + E'$

$ $

回顾 RLWE 密码系统（\autoref{sec:rlwe}）的另一种版本（\autoref{subsec:glwe-alternative}），明文与密文对满足以下关系：

$\Delta M(X) + E(X) = A(X)\cdot S(X) + B(X) \approx \Delta M(X)$

$ $

在上述关系中应用 $X = X^{J(h)}$，可推导出以下关系：

$\Delta M(X^{J(h)}) + E(X^{J(h)}) = A(X^{J(h)}) \cdot S(X^{J(h)}) + B(X^{J(h)}) \approx \Delta M(X^{J(h)})$

$ $

这一关系意味着，以 $S(X^{J(h)})$ 为秘密密钥解密密文 $C^{\langle h \rangle} = (A(X^{J(h)}), B(X^{J(h)}))$，即可得到 $\Delta M(X^{J(h)})$。因此，$C^{\langle h \rangle} = (A(X^{J(h)}), B(X^{J(h)}))$ 就是我们要寻找的 RLWE 密文，因为对它解密，再对其明文 $M(X^{J(h)})$ 解码，就能得到输入向量 $\vec{v}^{\langle h \rangle}$。

逐分量应用自同构，得到密文 $C^{\langle h\rangle}=(A(X^{J(h)}),B(X^{J(h)}))$，其秘密密钥为 $S(X^{J(h)})$。随后，使用 \autoref{subsec:bfv-key-switching} 中的小结~\ref*{subsec:bfv-key-switching} 将密钥切换回原密钥 $S(X)$。这会产生新的密文分量，并增加密钥切换误差；未经改变的同一分量对并不同时对应两个秘密密钥。下式只是在明文记号中省略该误差：

$ $

\noindent $\underbrace{\textsf{RLWE}_{S(X),\sigma}\bm{(}\Delta M(X^{J(h)})\bm{)}}_{\substack{\text{以下运算的结果：} \\ \text{明文与密文相加}}} = \underbrace{\bm{(} \text{ } 0, B(X^{J(h)}) \text{ } \bm{)}}_{\substack{\text{明文 } B(X^{J(h)}) \\ \text{（平凡密文）}}} + \bm{\langle} \text{ } \underbrace{\textsf{Decomp}^{\beta, l}\bm{(}A(X^{J(h)})\bm{)}, \text{ } \textsf{RLev}_{S(X), \sigma}^{\beta, l}\bm{(}S(X^{J(h)})\bm{)} \text{ } \bm{\rangle}}_{\substack{\substack{\text{加密 $A(X^{J(h)}) \cdot S(X^{J(h)})$ 的 RLWE 密文}\\ \text{其密钥已按 $S(X^{J(h)})\rightarrow S(X)$ 切换}}}}$



\subsubsection{小结}
\label{subsubsec:bfv-rotation-summary}

BFV 输入向量的旋转过程总结如下：

\begin{tcolorbox}[title={\textbf{\tboxlabel{\ref*{subsec:bfv-rotation}} BFV 输入向量槽位的同态旋转}}]

为支持输入向量槽位旋转，将 \textsf{Encoding\textsubscript{1}} 中的原始编码矩阵更新如下：

$ $

{\footnotesize{\noindent $\hathat{W} = \begin{bmatrix}
1 & 1 & \cdots & 1 & 1 & 1 & \cdots & 1\\
(\omega^{J(\frac{n}{2} - 1)}) & (\omega^{J(\frac{n}{2} - 2)}) & \cdots & (\omega^{J(0)}) & (\omega^{J_*(\frac{n}{2} - 1)}) & (\omega^{J_*(\frac{n}{2} - 2)}) & \cdots & (\omega^{J_*(0)})\\
(\omega^{J(\frac{n}{2} - 1)})^2 & (\omega^{J(\frac{n}{2} - 2)})^2 & \cdots & (\omega^{J(0)})^2 & (\omega^{J_*(\frac{n}{2} - 1)})^2 & (\omega^{J_*(\frac{n}{2} - 2)})^2 & \cdots & (\omega^{J_*(0)})^2 \\
\vdots & \vdots & \ddots & \vdots & \vdots & \vdots & \ddots & \vdots \\
(\omega^{J(\frac{n}{2} - 1)})^{n-1} & (\omega^{J(\frac{n}{2} - 2)})^{n-1} & \cdots & (\omega^{J(0)})^{n-1} & (\omega^{J_*(\frac{n}{2} - 1)})^{n-1} & (\omega^{J_*(\frac{n}{2} - 2)})^{n-1} & \cdots  & (\omega^{J_*(0)})^{n-1}
\end{bmatrix}$}}

$ $

$ $


，并将 \textsf{Decoding\textsubscript{2}} 中的原始解码矩阵更新如下：

$\hathat{W}^* = \begin{bmatrix}
1 & (\omega^{J(0)}) & (\omega^{J(0)})^2 & \cdots & (\omega^{J(0)})^{n-1}\\
1 & (\omega^{J(1)}) & (\omega^{J(1)})^2 & \cdots & (\omega^{J(1)})^{n-1}\\
1 & (\omega^{J(2)}) & (\omega^{J(2)})^2 & \cdots & (\omega^{J(2)})^{n-1}\\
\vdots & \vdots & \vdots & \ddots & \vdots \\
1 & (\omega^{J(\frac{n}{2}-1)}) & (\omega^{J(\frac{n}{2}-1)})^2 & \cdots & (\omega^{J(\frac{n}{2}-1)})^{n-1}\\
1 & (\omega^{J_*(0)}) & (\omega^{J_*(0)})^2 & \cdots & (\omega^{J_*(0)})^{n-1}\\
1 & (\omega^{J_*(1)}) & (\omega^{J_*(1)})^2 & \cdots & (\omega^{J_*(1)})^{n-1}\\
1 & (\omega^{J_*(2)}) & (\omega^{J_*(2)})^2 & \cdots & (\omega^{J_*(2)})^{n-1}\\
\vdots & \vdots & \vdots & \ddots & \vdots \\
1 & (\omega^{J_*(\frac{n}{2}-1)}) & (\omega^{J_*(\frac{n}{2}-1)})^2 & \cdots & (\omega^{J_*(\frac{n}{2}-1)})^{n-1}\\
\end{bmatrix}$

$ $

$ $

，其中，$J(h)$ 是\textit{旋转辅助公式}：$J(h) = 5^h \bmod 2n$，
\text{ } $J_*(h) = -5^h \bmod 2n$. 

$ $

使用 $\hathat{W}$ 进行编码：$\vec{m} = n^{-1}\cdot \hathat{W} \cdot I_n^R \cdot \vec{v}$。使用 $\hathat{W}^*$ 进行解码：$\vec{v} = \hathat{W}^* \cdot \vec{m}$

$ $

假设有如下 RLWE 密文和密钥切换密钥：

$\textsf{RLWE}_{S, \sigma}(\Delta M) = (A, B)$, \text{ } $\textsf{RLev}_{S, \sigma}^{\beta, l}\bm(S(X^{J(h)})\bm)$

$ $

那么，要使密文所对应原始输入向量 $\vec{v}$ 的前半部分元素在其内部向左循环旋转 $h$ 个位置，同时使 $\vec{v}$ 的后半部分元素也在其内部向左循环旋转 $h$ 个位置，过程如下：

\begin{enumerate}
\item 将 $A(X)$、$B(X)$ 更新为 $A(X^{J(h)})$、$B(X^{J(h)})$。
\item 执行如下从 $S(X^{J(h)})$ 到 $S(X)$ 的密钥切换（\autoref{subsec:bfv-key-switching}）：

$\textsf{RLWE}_{S(X),\sigma}\bm{(}\Delta M(X^{J(h)})\bm{)} = \bm{(} 0, B(X^{J(h)}) \bm{)} \text{ } + \text{ } \bm{\langle}  \textsf{Decomp}^{\beta, l}\bm{(}A(X^{J(h)})\bm{)}, \text{ } \textsf{RLev}_{S(X), \sigma}^{\beta, l}\bm{(}S(X^{J(h)})\bm{)} \bm{\rangle}$
\end{enumerate}


\end{tcolorbox}


\subsubsection{编码示例}
\label{subsubsec:bfv-batch-encoding-ex}

%$bfv_vector = [10, 0, 5, 13]
%bfv_vector2 = [2, 4, 8, 6]


在上述示例中，计算 $\bmod t$ 和 $\bmod q$ 时采用无符号模表示（例如 $[0, t-1], [0, q]$）。不过，即使使用居中模表示（例如 $[-\dfrac{t}{2}, \dfrac{t}{2})$、$[-\dfrac{q}{2}, \dfrac{q}{2})$），结果的正确性仍然成立。在许多实际的全同态加密库中，通常使用居中模运算来更有效地管理噪声增长。

假设采用如下设置：

$\mu = 8, \, n = \dfrac{\mu}{2} = 4, \, t = 17, \, q = 2^8 = 256, \, n^{-1} = 13, \, \Delta = \left\lfloor\dfrac{q}{t}\right\rfloor = 15$

$ $

$\mathcal{R}_{\langle 4, 17 \rangle} = \mathbb{Z}_{17}[X] / (X^4 + 1)$

$ $

$X^4 + 1 \pmod{17}$ 的根为 $X = \{2, 8, 15, 9\}$，验证如下：

$2^4 \equiv 8^4 \equiv 15^4 \equiv 9^4 \equiv 16 \equiv -1 \bmod{17}$

$ $

定义~\ref*{subsec:cyclotomic-def}（位于 \autoref{subsec:cyclotomic-def}）指出，第 $\mu$ 个分圆多项式的根是本原 $\mu$ 次单位根。而定义~\ref*{subsec:roots-def}（位于 \autoref{subsec:roots-def}）指出，本原 $\mu$ 次单位根的阶为 $\mu$。这些定义同时适用于 $X\in\mathbb{C}$（复数）上的分圆多项式和 $X \in \mathbb{Z}_t$（环）上的分圆多项式。

由于 $\{2, 8, 15, 9\}$ 是环 $\mathbb{Z}_{17}$ 上第 $(\mu=8)$ 个分圆多项式 $X^4 + 1$ 的根，因此它们也是 $\mathbb{Z}_{17}$ 的本原 $(\mu=8)$ 次单位根。所以，它们的阶（\autoref{subsec:order-def}）为 $\mu=8$，验证如下：

$2^8 \equiv 8^8 \equiv 15^8 \equiv 9^8 \equiv 1 \bmod 17$

$2^4 \equiv 8^4 \equiv 15^4 \equiv 9^4 \equiv 16 \not\equiv 1 \bmod 17$

$ $

定义~\ref*{subsec:cyclotomic-def}（位于 \autoref{subsec:cyclotomic-def}）和定理~\ref*{subsec:cyclotomic-theorem}（\autoref{subsec:cyclotomic-theorem}）还指出，对于每个本原 $\mu$ 次单位根 $\omega$，$\{\omega^k\}_{\textsf{gcd}(k, \mu) = 1}$ 生成第 $\mu$ 个分圆多项式的全部根。注意，对于第 $(\mu=8)$ 个分圆多项式 $X^4 + 1$，它的根 $\{2, 8, 15, 9\}$ 按如下方式生成全部 $(\mu=8)$ 次单位根：

$\{2^1, 2^3, 2^5, 2^7\} \equiv \{8^1, 8^3, 8^5, 8^7\} \equiv \{15^1, 15^3, 15^5, 15^7\} \equiv \{9^1, 9^3, 9^5, 9^7\} \equiv \{2, 8, 15, 9\} \bmod 17$

$ $

在 $X^4 + 1$ 的根 $\{2, 8, 15, 9\}$ 中，选取 $\omega = 9$ 作为基础根，按如下方式构造编码矩阵 $\hathat{W}$ 和解码矩阵 $\hathat{W}^*$：

$\hathat{W} = \begin{bmatrix}
1 & 1 & 1 & 1\\
\omega^{J(1)} & \omega^{J(0)} & \omega^{J_*(1)} & \omega^{J_*(0)}\\
(\omega^{J(1)})^2 & (\omega^{J(0)})^2 & (\omega^{J_*(1)})^2 & (\omega^{J_*(0)})^2\\
(\omega^{J(1)})^3 & (\omega^{J(0)})^3 & (\omega^{J_*(1)})^3 & (\omega^{J_*(0)})^3\\
\end{bmatrix}$ \textcolor{red}{  $\rhd$ 其中 $J(h) = 5^h \bmod 8$}

$ = \begin{bmatrix}
1 & 1 & 1 & 1\\
9^{5} & 9^{1} & 9^{3} & 9^{7}\\
(9^{5})^2 & (9^{1})^2 & (9^{3})^2 & (9^{7})^2\\
(9^{5})^3 & (9^{1})^3 & (9^{3})^3 & (9^{7})^3\\
\end{bmatrix} \equiv \begin{bmatrix}
1 & 1 & 1 & 1\\
8 & 9 & 15 & 2\\
13 & 13 & 4 & 4\\
2 & 15 & 9 & 8\\
\end{bmatrix} \bmod{17}$ 

$ $ 

$\hathat{W}^* = \begin{bmatrix}
1 & \omega^{J(0)} & (\omega^{J(0)})^2 & (\omega^{J(0)})^3\\
1 & \omega^{J(1)} & (\omega^{J(1)})^2 & (\omega^{J(1)})^3\\
1 & \omega^{J_*(0)} & (\omega^{J_*(0)})^2 & (\omega^{J_*(0)})^3\\
1 & \omega^{J_*(1)} & (\omega^{J_*(1)})^2 & (\omega^{J_*(1)})^3\\
\end{bmatrix} \equiv \begin{bmatrix}
1 & 9 & 13 & 15\\
1 & 8 & 13 & 2\\
1 & 2 & 4 & 8\\
1 & 15 & 4 & 9\\
\end{bmatrix}  \bmod{17}$

$ $

注意，定理~\ref*{subsec:vandermonde-euler-integer-ring}（位于 \autoref{subsec:vandermonde-euler-integer-ring}）可以验证如下：

$\hathat{W}^* \cdot \hathat{W} = \begin{bmatrix}
1 & 9 & 13 & 15\\
1 & 8 & 13 & 2\\
1 & 2 & 4 & 8\\
1 & 15 & 4 & 9\\
\end{bmatrix} \cdot \begin{bmatrix}
1 & 1 & 1 & 1\\
8 & 9 & 15 & 2\\
13 & 13 & 4 & 4\\
2 & 15 & 9 & 8\\
\end{bmatrix} = \begin{bmatrix}
0 & 0 & 0 & 4\\
0 & 0 & 4 & 0\\
0 & 4 & 0 & 0\\
4 & 0 & 0 & 0\\
\end{bmatrix} = n I_n^{R} \pmod{17}$

$ $

现在，假设要对 $\mathbb{Z}_{17}$ 中的以下两个输入向量（即输入向量槽位）进行编码：

$\vec{v}_1 = (10, 3, 5, 13)$

$\vec{v}_2 = (2, 4, 3, 6)$

$\vec{v}_1 + \vec{v}_2 = (10, 3, 5, 13) + (2, 4, 3, 6) \equiv (12, 7, 8, 2) \bmod 17$

$ $

这两个向量的编码如下：

$ $

$\vec{m}_1 = n^{-1} \hathat{W} \cdot I_n^R \cdot \vec{v} = 13 \cdot \begin{bmatrix}
1 & 1 & 1 & 1\\
8 & 9 & 15 & 2\\
13 & 13 & 4 & 4\\
2 & 15 & 9 & 8\\
\end{bmatrix} \cdot \begin{bmatrix}
0 & 0 & 0 & 1\\
0 & 0 & 1 & 0\\
0 & 1 & 0 & 0\\
1 & 0 & 0 & 0\\
\end{bmatrix} \cdot \begin{bmatrix}
10\\
3\\
5\\
13\\
\end{bmatrix} \equiv (12, 11, 12, 1)  \bmod 17$

$ $

$\vec{m}_2 = n^{-1} \hathat{W} \cdot I_n^R \cdot \vec{v} = 13 \cdot \begin{bmatrix}
1 & 1 & 1 & 1\\
8 & 9 & 15 & 2\\
13 & 13 & 4 & 4\\
2 & 15 & 9 & 8\\
\end{bmatrix} \cdot \begin{bmatrix}
0 & 0 & 0 & 1\\
0 & 0 & 1 & 0\\
0 & 1 & 0 & 0\\
1 & 0 & 0 & 0\\
\end{bmatrix} \cdot \begin{bmatrix}
2\\
4\\
3\\
6\\
\end{bmatrix} \equiv (8, 5, 14, 6) \bmod 17$

$ $

$ $

$\Delta M_1(X) = 15\cdot(12 + 11X + 12X^2 + 1X^3) = 180 + 165X + 180X^2 + 15X^3 \pmod{q}$ \textcolor{red}{ $\rhd$ 其中 $q = 256$}

$\Delta M_2(X) = 15\cdot(8 + 5X + 14X^2 + 6X^3) = 120 + 75X + 210X^2 + 90X^3 \pmod{q}$

$\Delta M_{1+2}(X) = \Delta M_1(X) + \Delta M_2(X) = \Delta (M_1(X) + M_2(X)) = 44 + 240X + 134X^2 + 105X^3 \pmod{q}$

$M_{1+2}(X) =\left\lceil\dfrac{17}{256}(44 + 240X + 134X^2 + 105X^3)\right\rfloor \bmod 17$

$= \left\lceil2.922 + 15.938X + 8.898X^2 + 6.973X^3\right\rfloor \bmod 17$

$3 + 16X + 9X^2 + 7X^3 \bmod 17 $

$ $


这一多项式与 $\vec{m}_{1+2}$ 的值一致，如下所示：

$\vec{m}_1 + \vec{m}_2 = (12, 11, 12, 1) + (8, 5, 14, 6) = (20, 16, 26, 7) \equiv (3, 16, 9, 7) \bmod 17$

$ $

最后，按如下方式解码 $\vec{m}_{1+2}$：



$\vec{v}_{1+2} = \hathat{W}^* \cdot \vec{m}_{1+2} = \begin{bmatrix}
1 & 9 & 13 & 15\\
1 & 8 & 13 & 2\\
1 & 2 & 4 & 8\\
1 & 15 & 4 & 9\\
\end{bmatrix} \cdot \begin{bmatrix}3\\16\\9\\7\end{bmatrix} = (12, 7, 8, 2) \bmod{17}$

$ $

这一结果与预期的向量和 $\vec{v}_1 + \vec{v}_2$ 一致。




\subsubsection{旋转示例}
\label{subsubsec:bfv-rotation-ex}

同样，在上述示例中，计算 $\bmod t$ 和 $\bmod q$ 时使用无符号模表示（例如 $[0, t-1], [0, q]$）。不过，即使使用居中模表示（例如 $[-\dfrac{t}{2}, \dfrac{t}{2})$、$[-\dfrac{q}{2}, \dfrac{q}{2})$），结果的正确性仍然成立。

假设采用如下设置：

$\mu = 16, n = \dfrac{\mu}{2} = 8, \text{ } t = 17, \text{ } q = 2^8 = 256, \text{ } n^{-1} = 15, \text{ } \Delta = 15$

$ $

$\mathcal{R}_{\langle 8, 17 \rangle} = \mathbb{Z}_{17}[X] / (X^8 + 1)$

$ $

$X^8 + 1 \pmod{17}$ 的根为 $X = \{3, 5, 6, 7, 10, 11, 12, 14\}$，验证如下：

$3^8 \equiv 5^8 \equiv 6^8 \equiv 7^8 \equiv 10^8 \equiv 11^8 \equiv 12^8 \equiv 14^8 \equiv 16 \equiv -1 \bmod{17}$


$ $

在 $X^8 + 1$ 的根 $\{3, 5, 6, 7, 10, 11, 12, 14\}$ 中，选取 $\omega = 3$ 作为基础根，按如下方式构造编码矩阵 $\hathat{W}$ 和解码矩阵 $\hathat{W}^*$：

$\hathat{W} = \begin{bmatrix}
1 & 1 & 1 & 1 & 1 & 1 & 1 & 1\\
\omega^{J(3)} & \omega^{J(2)} & \omega^{J(1)} & \omega^{J(0)} & \omega^{J_*(3)} & \omega^{J_*(2)} & \omega^{J_*(1)} & \omega^{J_*(0)}\\
(\omega^{J(3)})^2 & (\omega^{J(2)})^2 & (\omega^{J(1)})^2 & (\omega^{J(0)})^2 & (\omega^{J_*(3)})^2 & (\omega^{J_*(2)})^2 & (\omega^{J_*(1)})^2 & (\omega^{J_*(0)})^2\\
(\omega^{J(3)})^3 & (\omega^{J(2)})^3 & (\omega^{J(1)})^3 & (\omega^{J(0)})^3 & (\omega^{J_*(3)})^3 & (\omega^{J_*(2)})^3 & (\omega^{J_*(1)})^3 & (\omega^{J_*(0)})^3\\
(\omega^{J(3)})^4 & (\omega^{J(2)})^4 & (\omega^{J(1)})^4 & (\omega^{J(0)})^4 & (\omega^{J_*(3)})^4 & (\omega^{J_*(2)})^4 & (\omega^{J_*(1)})^4 & (\omega^{J_*(0)})^4\\
(\omega^{J(3)})^5 & (\omega^{J(2)})^5 & (\omega^{J(1)})^5 & (\omega^{J(0)})^5 & (\omega^{J_*(3)})^5 & (\omega^{J_*(2)})^5 & (\omega^{J_*(1)})^5 & (\omega^{J_*(0)})^5\\
(\omega^{J(3)})^6 & (\omega^{J(2)})^6 & (\omega^{J(1)})^6 & (\omega^{J(0)})^6 & (\omega^{J_*(3)})^6 & (\omega^{J_*(2)})^6 & (\omega^{J_*(1)})^6 & (\omega^{J_*(0)})^6\\
(\omega^{J(3)})^7 & (\omega^{J(2)})^7 & (\omega^{J(1)})^7 & (\omega^{J(0)})^7 & (\omega^{J_*(3)})^7 & (\omega^{J_*(2)})^7 & (\omega^{J_*(1)})^7 & (\omega^{J_*(0)})^7\\
\end{bmatrix}$ 

$ \equiv \begin{bmatrix}
1 & 1 & 1 & 1 & 1 & 1 & 1 & 1\\
12&14&5&3&10&11&7&6\\
8&9&8&9&15&2&15&2\\
11&7&6&10&14&5&3&12\\
13&13&13&13&4&4&4&4\\
3&12&14&5&6&10&11&7\\
2&15&2&15&9&8&9&8\\
7&6&10&11&5&3&12&14\\
\end{bmatrix} \bmod{17}$ 

$ $ 

$\hathat{W}^* = \begin{bmatrix}
1 & \omega^{J(0)} & (\omega^{J(0)})^2 & (\omega^{J(0)})^3 & (\omega^{J(0)})^4 & (\omega^{J(0)})^5 & (\omega^{J(0)})^6 & (\omega^{J(0)})^7\\
1 & \omega^{J(1)} & (\omega^{J(1)})^2 & (\omega^{J(1)})^3 & (\omega^{J(1)})^4& (\omega^{J(1)})^5& (\omega^{J(1)})^6& (\omega^{J(1)})^7\\
1 & \omega^{J(2)} & (\omega^{J(2)})^2 & (\omega^{J(2)})^3 & (\omega^{J(2)})^4 & (\omega^{J(2)})^5 & (\omega^{J(2)})^6 & (\omega^{J(2)})^7\\
1 & \omega^{J(3)} & (\omega^{J(3)})^2 & (\omega^{J(3)})^3 & (\omega^{J(3)})^4& (\omega^{J(3)})^5& (\omega^{J(3)})^6& (\omega^{J(3)})^7\\
1 & \omega^{J_*(0)} & (\omega^{J_*(0)})^2 & (\omega^{J_*(0)})^3 & (\omega^{J_*(0)})^4 & (\omega^{J_*(0)})^5 & (\omega^{J_*(0)})^6 & (\omega^{J_*(0)})^7\\
1 & \omega^{J_*(1)} & (\omega^{J_*(1)})^2 & (\omega^{J_*(1)})^3 & (\omega^{J_*(1)})^4& (\omega^{J_*(1)})^5& (\omega^{J_*(1)})^6& (\omega^{J_*(1)})^7\\
1 & \omega^{J_*(2)} & (\omega^{J_*(2)})^2 & (\omega^{J_*(2)})^3 & (\omega^{J_*(2)})^4 & (\omega^{J_*(2)})^5 & (\omega^{J_*(2)})^6 & (\omega^{J_*(2)})^7\\
1 & \omega^{J_*(3)} & (\omega^{J_*(3)})^2 & (\omega^{J_*(3)})^3 & (\omega^{J_*(3)})^4& (\omega^{J_*(3)})^5& (\omega^{J_*(3)})^6& (\omega^{J_*(3)})^7\\
\end{bmatrix}$

$ \equiv \begin{bmatrix}
1&3&9&10&13&5&15&11\\
1&5&8&6&13&14&2&10\\
1&14&9&7&13&12&15&6\\
1&12&8&11&13&3&2&7\\
1&6&2&12&4&7&8&14\\
1&7&15&3&4&11&9&12\\
1&11&2&5&4&10&8&3\\
1&10&15&14&4&6&9&5\\
\end{bmatrix}  \bmod{17}$

$ $

$ $

注意，定理~\ref*{subsec:vandermonde-euler-integer-ring}（位于 \autoref{subsec:vandermonde-euler-integer-ring}）可以验证如下：

$\hathat{W}^* \cdot \hathat{W} = \begin{bmatrix}
1&3&9&10&13&5&15&11\\
1&5&8&6&13&14&2&10\\
1&14&9&7&13&12&15&6\\
1&12&8&11&13&3&2&7\\
1&6&2&12&4&7&8&14\\
1&7&15&3&4&11&9&12\\
1&11&2&5&4&10&8&3\\
1&10&15&14&4&6&9&5\\
\end{bmatrix} \cdot \begin{bmatrix}
1 & 1 & 1 & 1 & 1 & 1 & 1 & 1\\
12&14&5&3&10&11&7&6\\
8&9&8&9&15&2&15&2\\
11&7&6&10&14&5&3&12\\
13&13&13&13&4&4&4&4\\
3&12&14&5&6&10&11&7\\
2&15&2&15&9&8&9&8\\
7&6&10&11&5&3&12&14\\
\end{bmatrix}$

$ = \begin{bmatrix}
0 & 0 & 0 & 0 & 0 & 0 & 0 & 8\\
0 & 0 & 0 & 0 & 0 & 0 & 8 & 0\\
0 & 0 & 0 & 0 & 0 & 8 & 0 & 0\\
0 & 0 & 0 & 0 & 8 & 0 & 0 & 0\\
0 & 0 & 0 & 8 & 0 & 0 & 0 & 0\\
0 & 0 & 8 & 0 & 0 & 0 & 0 & 0\\
0 & 8 & 0 & 0 & 0 & 0 & 0 & 0\\
8 & 0 & 0 & 0 & 0 & 0 & 0 & 0\\
\end{bmatrix} = n I_n^{R} \pmod{17}$

$ $

现在，假设要对 $\mathbb{Z}_{17}$ 中的以下输入向量（即输入向量槽位）进行编码：

$\vec{v} = (1, 2, 3, 4, 5, 6, 7, 8)$

将这一向量向左旋转 3 个位置（即前半部分槽位与后半部分槽位分别在各自的组内循环回绕），得到一个新向量：

$\vec{v}_r = (4, 1, 2, 3, 8, 5, 6, 7) \bmod 17$

$ $

$\vec{v}$ 的编码如下：

$ $

$\vec{m} = n^{-1} \hathat{W} \cdot I_n^R \cdot \vec{v} $

$= 15 \cdot  \begin{bmatrix}
1 & 1 & 1 & 1 & 1 & 1 & 1 & 1\\
12&14&5&3&10&11&7&6\\
8&9&8&9&15&2&15&2\\
11&7&6&10&14&5&3&12\\
13&13&13&13&4&4&4&4\\
3&12&14&5&6&10&11&7\\
2&15&2&15&9&8&9&8\\
7&6&10&11&5&3&12&14\\
\end{bmatrix} \cdot \begin{bmatrix}
0 & 0 & 0 & 0 & 0 & 0 & 0 & 1\\
0 & 0 & 0 & 0 & 0 & 0 & 1 & 0\\
0 & 0 & 0 & 0 & 0 & 1 & 0 & 0\\
0 & 0 & 0 & 0 & 1 & 0 & 0 & 0\\
0 & 0 & 0 & 1 & 0 & 0 & 0 & 0\\
0 & 0 & 1 & 0 & 0 & 0 & 0 & 0\\
0 & 1 & 0 & 0 & 0 & 0 & 0 & 0\\
1 & 0 & 0 & 0 & 0 & 0 & 0 & 0\\
\end{bmatrix} \cdot \begin{bmatrix}
1\\
2\\
3\\
4\\
5\\
6\\
7\\
8\\
\end{bmatrix}$

$ \equiv (13, 16, 10, 5, 9, 12, 7, 1)  \bmod 17$


$ $

$\Delta M(X) = 15\cdot(13 + 16X + 10X^2 + 5X^3 + 9X^4 + 12X^5 + 7X^6 + X^7)$


$ = 195 + 240X + 150X^2 + 75X^3 + 135X^4 + 180X^5 + 105X^6 + 15X^7 \pmod{q}$ \textcolor{red}{ $\rhd$ 其中 $q = 256$}

$ $

这一多项式与 $\Delta M(X^{J(3)})$ 的值一致，如下所示：

$\Delta M(X^{J(3)}) = \Delta M(X^{13}) $

$= 195 + 180X - 150X^2 - 15X^3 + 135X^4 - 240X^5 - 105X^6 + 75X^7 \pmod{q}$

$ $



%Note that the coefficients of the scaled polynomials $\Delta M(X^{J(3)})$ are still within the range of the ciphertext domain $q=64$ (which must hold throughout all homomorphic computations to preserve correctness). 

现在，按如下方式解码 $M(X^{J(3)})$：

$\vec{m}_{J(3)} = \dfrac{\Delta \vec{m}_{J(3)}}{\Delta} = \dfrac{(195, 180, -150, -15, 135, -240, -105, 75)}{15} = (13, 12, -10, -1, 9, -16, -7, 5) \pmod{17}$

$ $


$\vec{v}_{J(3)} = \hathat{W}^* \cdot \vec{m}_{J(3)} = \begin{bmatrix}
1&3&9&10&13&5&15&11\\
1&5&8&6&13&14&2&10\\
1&14&9&7&13&12&15&6\\
1&12&8&11&13&3&2&7\\
1&6&2&12&4&7&8&14\\
1&7&15&3&4&11&9&12\\
1&11&2&5&4&10&8&3\\
1&10&15&14&4&6&9&5\\
\end{bmatrix} \cdot \begin{bmatrix}
13\\
12\\
-10\\
-1\\
9\\
-16\\
-7\\
5\end{bmatrix}$

$= (4, 1, 2, 3, 8, 5, 6, 7) \pmod{17}$

$= \vec{v}_r$

$ $

解码得到的 $\vec{v}_{J(3)}$ 与预期的旋转后输入向量 $\vec{v}_r$ 一致。


$ $

实际中，对加密分量应用自同构，得到
\[
\textsf{ct}^{\langle h=3\rangle}=(A(X^{J(3)}),B(X^{J(3)}))
\]
。其解密秘密密钥为 $S(X^{J(3)})$。在执行通常的解密之前，使用相应旋转密钥将秘密密钥切换回 $S(X)$；只要变换后的误差与密钥切换误差仍满足界限，解码就得到 $\vec v_r$。

$ $




\para{源代码：} 运行 \href{https://github.com/wang2000star/fhe-textbook-cn/blob/chinese-edition/source%20code/bfv.py}{\underline{这个 Python 脚本}}（BFV 加法）和 \href{https://github.com/wang2000star/fhe-textbook-cn/blob/chinese-edition/source%20code/bfv_rotation_only.py}{\underline{这个脚本}}（仅 BFV 旋转），可以执行 BFV 批量编码和同态输入向量旋转的示例。



\subsection{应用：矩阵乘法}
\label{subsec:bfv-matrix-multiplication}

BFV 可以先逐槽相乘，再通过旋转和加法求和，计算打包内积。在槽位数为二次幂的单个循环块内，倍增求和只需对数数量的旋转；若要合并两行，还需要行交换。下述矩阵与向量乘法的对角方法假设 $m\leq n$，且 $\rho$ 表示对全部 $n$ 个有效槽位的逻辑循环旋转。前节只实现两个半行各自旋转，因此仅凭该操作不能实现完整的 $\rho$。可以将逻辑向量放在更大环的一个循环行中，也可组合行内旋转、行交换自同构 $X\mapsto X^{-1}$ 和明文掩码。给定矩阵 $A$ 和向量 $\vec x$（\autoref{subsec:matrix-arithmetic} 中的定义~\ref*{subsec:matrix-arithmetic}）：

$A = \begin{bmatrix}
 a_{\langle 0, 0\rangle} & a_{\langle 0, 1\rangle} & a_{\langle 0, 2\rangle} & \cdots & a_{\langle 0, n-1\rangle}\\
 a_{\langle 1, 0\rangle} & a_{\langle 1, 1\rangle} & a_{\langle 1, 2\rangle} & \cdots & a_{\langle 1, n-1\rangle} \\
 a_{\langle 2, 0\rangle} & a_{\langle 2, 1\rangle} & a_{\langle 2, 2\rangle} & \cdots & a_{\langle 2, n-1\rangle} \\
\vdots & \vdots & \vdots & \ddots & \vdots \\
 a_{\langle m-1, 0\rangle} & a_{\langle m-1, 1\rangle} & a_{\langle m-1, 2\rangle} & \cdots & a_{\langle m-1, n-1\rangle} \\
\end{bmatrix} = \begin{bmatrix} 
\vec{a}_{\langle 0, * \rangle} \\ 
\vec{a}_{\langle 1, * \rangle} \\ 
\vec{a}_{\langle 2, * \rangle} \\ 
\vdots\\
\vec{a}_{\langle m-1, * \rangle} \\ 
\end{bmatrix}, \text{ } \vec{x} = (x_0, x_1, \cdots, x_{n-1})$

$ $

$A \cdot \vec{x}$ 的结果是一个 $m$ 维向量，计算方式如下：

$A \cdot \vec{x} = \Big(\vec{a}_{\langle 0, * \rangle} \cdot \vec{x}, \text{ } \vec{a}_{\langle 1, * \rangle} \cdot \vec{x}, \text{ }\cdots,\text{ } \vec{a}_{\langle m-1, * \rangle} \cdot \vec{x} \Big) = \left(\sum\limits_{i=0}^{n-1} a_{0, i} \cdot  x_i, \text{ } \sum\limits_{i=0}^{n-1}  a_{1, i} \cdot x_i, \cdots, \text{ } \sum\limits_{i=0}^{n-1} a_{m-1, i} \cdot x_i \right)$

$ $

将 $\vec{v}$ 向左旋转 $h$ 个位置后的结果定义为 $\rho(\vec{v}, h)$。回顾 Hadamard 乘积（\autoref{subsec:vector-arithmetic} 中的定义~\ref*{subsec:vector-arithmetic}），它定义为两个向量逐槽位相乘：

$\vec{a} \odot \vec{b} = (a_0 b_0, \text{ } a_1 b_1, \text{ } \cdots, \text{ } a_{n-1} b_{n-1})$

$ $

从矩阵 $A$ 中提取 $n$ 个不同的对角向量 $\vec{u}_i$，定义如下：

$\vec{u}_i = \{a_{\langle j \bmod m, \text{ } (i + j)\bmod n \rangle}\}_{j = 0}^{n-1}$

$ $

那么，原始矩阵与向量相乘的公式可以等价地构造如下：

$A\cdot\vec x=\left(\sum_{i=0}^{n-1}\vec u_i\odot\rho(\vec x,i)\right)_{0:m}$

，其计算结果等价于 $A \cdot \vec{x}$。
上述公式可以用于同态计算，因为 BFV 支持通过输入向量的多项式编码形式之间的密文与明文相乘，计算输入向量的 Hadamard 乘积（\autoref{subsec:bfv-mult-plain}）；BFV 也支持将 $\rho(\vec{v}, h)$ 作为同态输入向量槽位旋转来计算（\autoref{subsec:bfv-rotation}）。在同态计算上述公式后，只需使用所得输入向量的前 $m$ 个槽位（共 $n$ 个槽位）来存储 $A \cdot \vec{x}$ 的结果。



\subsection{噪声自举}
\label{subsec:bfv-bootstrapping}

\noindent \textbf{- 参考资料 1：}
\href{https://eprint.iacr.org/2022/1363.pdf}{重新审视 BGV 与 BFV 的自举}~\cite{cryptoeprint:2022/1363}


\noindent \textbf{- 参考资料 2：}
\href{https://eprint.iacr.org/2014/873.pdf}{HELib 的自举}~\cite{10.1007/s00145-020-09368-7}

\noindent \textbf{- 参考资料 3：}
\href{https://arxiv.org/pdf/1906.02867}{关于自举中低位提取多项式的说明}~\cite{huo2019notelowerdigitsextraction}

\noindent \textbf{- 参考资料 4：}
\href{https://eprint.iacr.org/2024/1587.pdf}{采用分圆素数模数的全同态加密}~\cite{cryptoeprint:2024/1587}

%\noindent \textbf{- Reference 4:} 
%\href{https://s-space.snu.ac.kr/bitstream/10371/152893/1/000000155273.pdf}{Bootstrapping Methods for Homomorphic Encryption}

%\noindent \textbf{- Reference 5:} 
%\href{https://www.esat.kuleuven.be/cosic/publications/thesis-418.pdf}{Bootstrapping Algorithms for BGV and FV}

$ $

在 BFV 中，连续进行密文与密文相乘会使噪声以乘法方式增长；一旦噪声溢出到消息位，消息就会遭到破坏。自举是重置已增长噪声的过程。

\subsubsection{总体思路}
\label{subsubsec:bfv-bootstrapping-high-level}

令 $t=p^r$、$\varepsilon>r$ 且 $P=p^\varepsilon$。下面分解解密函数，给出一般系数与槽位变换构造，并在 \autoref{subsubsec:bfv-bootstrapping-summary} 给出显式的逐门正确性充分条件。该条件中的同态多项式求值采用标准 FV：先计算整数张量乘积，再按当前明文模数 $T$ 以 $T/q$ 缩放；当 $T\nmid q$ 时，不能改用前述 $1/\Delta$ 教学变体。这是条件性正确性结论，并不是安全或优化参数的选择。使用秘密密钥在自身之下的加密，还需所引构造的求值密钥安全及循环安全性假设~\cite{cryptoeprint:2022/1363}。

对于当前明文模数 $T$，按精确的有理尺度 $\mu_T=q/T$ 衡量误差会更方便。令 $r_T=q-T\lfloor q/T\rfloor$，同一个密文满足
\[
\phi_S(C)\equiv\mu_TM_T+\mathcal E_T\pmod q,
\qquad \mathcal E_T=E_T-\frac{r_T}{T}M_T.
\]
这里 $M_T$ 的各系数取中心代表。若 $\|\mathcal E_T\|_\infty<q/(2T)$，解码就正确。改变误差坐标既不改变密文，也不把方案变成近似明文方案；这个条件正是 \autoref{subsec:bfv-enc-dec} 的解码条件。

对于相位满足 $AS+B=\Delta M+E\pmod q$ 的密文 $(A,B)$，目标明文恒等式为
\[
A'=\Bigl[\Bigl\lceil(P/q)A\Bigr\rfloor\Bigr]_P,\quad
B'=\Bigl[\Bigl\lceil(P/q)B\Bigr\rfloor\Bigr]_P,\quad
Z=[A'S+B']_P,
\qquad M=\Bigl[\Bigl\lceil Z/p^{\varepsilon-r}\Bigr\rfloor\Bigr]_{p^r}.
\]
\begin{enumerate}
\item \textbf{模数切换：}将输入密文模数从 $q$ 改为 $P$。这一步本身不改变原明文模数 $p^r$。参数适当时，切换后相位为 $p^{\varepsilon-r}M+E'\pmod P$，其中 $\|E'\|_\infty<p^{\varepsilon-r}/2$。
\item \textbf{同态解密：}使用在较大密文模数 $q$ 下加密 $S$ 的求值密文，此时明文模数为 $P$，尺度为 $\Delta'=\lfloor q/P\rfloor$。计算 $A'\cdot\textsf{RLWE}_{S,\sigma}(\Delta'S)+(0,\Delta'B')$。结果加密相位 $Z\pmod P$，并含有额外求值噪声。公开常数 $B'$ 必须以尺度 $\Delta'$ 编码。
\item \textbf{\textsf{CoeffToSlot}：}应用同态线性变换，将 $Z$ 的系数放入槽位。在满足 \autoref{subsubsec:bfv-rotation-summary} 的完全标量分裂条件时，对槽位向量应用的变换为 $n^{-1}\hathat W I_n^R$。
\item \textbf{数字提取：}在每个槽位中，通过低位提取多项式及精确除法步骤，计算 $[\lceil z_i/p^{\varepsilon-r}\rfloor]_{p^r}$。当 $p$ 为奇数时使用平衡数字；当 $p=2$ 时，删除无符号数字实现的是向下取整，因此删除之前需要加半单位偏移量。下文将说明这些数字恒等式；它们用于完整自举时，还须满足 \autoref{subsubsec:bfv-bootstrapping-coefftoslot} 的 CRT 适用域限制。
\item \textbf{\textsf{SlotToCoeff}：}应用逆线性变换 $\hathat W^*$，将恢复出的系数放回多项式。结果加密 $M$，并带有求值产生的新误差 $E^{\langle b\rangle}$。其界取决于参数和电路，并不自动小于任意输入误差 $E$，也不能直接与尺度不同的 $E'$ 比较。
\end{enumerate}

旧的 $E'$ 在求值舍入函数时被消除，而同态求值自身产生的误差仍然保留。可用的自举必须为后续操作留下足够的输出噪声预算。只有选择适当参数并控制总失败概率，重复自举才能支持任意预先给定的有限电路深度。

\subsubsection{模数切换}
\label{subsubsec:bfv-bootstrapping-modulus-switch}

令 $P=p^\varepsilon$，在约减前选用相容的整数提升，定义各分量舍入多项式 $R_A=A'-(P/q)A$ 和 $R_B=B'-(P/q)B$。其系数绝对值至多为 $1/2$。切换后的精确相位为
\[
A'S+B'\equiv p^{\varepsilon-r}M+E'\pmod P,
\quad E'=\frac PqE+\left(\frac Pq\Delta-p^{\varepsilon-r}\right)M+R_AS+R_B.
\]
因此，目标舍入的一个充分条件是 $\|E'\|_\infty<p^{\varepsilon-r}/2$。各分量舍入项不能用仅对旧误差舍入来替代。除缩放后的旧误差和尺度失配项外，该项的确定性界为 $(\|S\|_1+1)/2$。

\subsubsection{同态解密}
\label{subsubsec:bfv-bootstrapping-homomorphic-decryption}

写成 $Z=[A'S+B']_P=[p^{\varepsilon-r}M+E']_P$。取 $\Delta'=\lfloor q/P\rfloor$，令求值密文 $(A_s,B_s)$ 满足
\[
A_sS+B_s=\Delta'S+E_s\pmod q.
\]
同态内积得到密文
\[
C=(A'A_s,\ A'B_s+\Delta'B'),
\qquad \phi_S(C)=\Delta'(A'S+B')+A'E_s\pmod q.
\]
这明确区分了加密明文 $\Delta'S$ 与加密秘密密钥 $S$。若 $A'S+B'=Z+PK$，则
\[
\phi_S(C)\equiv\Delta'Z+A'E_s-r_PK\pmod q,
\qquad r_P=q-P\Delta'.
\]
因此，当 $P\nmid q$ 时，明文模 $P$ 约减表现为额外误差 $-r_PK$，并不需要单独执行密文侧的取模门。另有 $\|A'E_s\|_\infty\leq\|A'\|_1\|E_s\|_\infty$。仅有宽泛的不等式 $q\gg P$，并不能保证该误差落在量级为 $q/P$ 的解码预算内。

目标系数操作是 $[\lceil z_i/p^{\varepsilon-r}\rfloor]_{p^r}$。它通过一系列多项式求值和精确除法实现，而不是由一个 $p$ 次多项式完成整个舍入函数。若要逐系数求值，先执行 \textsf{CoeffToSlot}，再执行 \autoref{subsubsec:bfv-bootstrapping-digit-extraction} 的数字提取过程，最后执行 \textsf{SlotToCoeff}。同态求值误差与加密整数 $z_i$ 内部的旧误差 $E'$ 是不同量。

\subsubsection{\textsf{CoeffToSlot} 与 \textsf{SlotToCoeff}}
\label{subsubsec:bfv-bootstrapping-coefftoslot}

下列 $n$ 个标量槽位的矩阵假设 $p\equiv1\pmod{2n}$，并将一个本原单位根 Hensel 提升到 $\mathbb Z_{p^\varepsilon}$，随后在每个明文模数下相容地约减。一般地，当 $\gcd(p,2n)=1$ 且 $d=\operatorname{ord}_{2n}(p)$ 时，有 $n/d$ 个扩域环槽位。下述迹对偶基构造将全部系数打包进 $d$ 个密文的标量槽位，因此不必完全分裂为标量槽位。这里的矩阵乘法表示由自同构、密钥切换和明文乘法组成的同态线性电路。这条 CRT 路线在本书二次幂环中仍要求 $p$ 为奇数：当 $p=2$ 时，$X^n+1=(X+1)^n\pmod2$ 含重因子。二进制 CRT 自举必须改用与 $2$ 互素的分圆指标，或采用另有独立论证的构造。

\textbf{一般可实现的线性变换。}为包含替代的二进制环，暂记 $R_j=(\mathbb Z/p^j\mathbb Z)[X]/\Phi_{m_c}(X)$，其中 $p\nmid m_c$、$n=\varphi(m_c)$、$d=\operatorname{ord}_{m_c}(p)$ 且 $\ell=n/d$。本章原来的环取 $m_c=2n$。选择模 $p^\varepsilon$ 下相容的 $\Phi_{m_c}$ 首一因子，再约减到模 $p^j$；CRT 给出 $\ell$ 个次数为 $d$ 的非分歧扩域环槽位。

对 $a\in(\mathbb Z/m_c\mathbb Z)^\times$，令 $\tau_a(X)=X^a$。由于模 $p$ 的代数是有限域的直积，迹配对非退化，其 Gram 行列式在每个模 $p^j$ 下都是单位。因此存在迹对偶基 $b_0,\ldots,b_{n-1}$，满足
\[
\operatorname{Tr}(b_uX^v)=\delta_{uv},\qquad
\operatorname{Tr}(W)=\sum_a\tau_a(W),\qquad
z_u=\operatorname{Tr}(b_uZ).
\]
于是，每个基环线性映射 $\mathcal T:R_j\to R_j$ 都有显式实现
\[
\mathcal T(Z)=\sum_a\kappa_a\tau_a(Z),\qquad
\kappa_a=\sum_{u=0}^{n-1}\mathcal T(X^u)\tau_a(b_u).
\]
代入迹的展开式，所得结果就是 $\sum_u\mathcal T(X^u)z_u$。每一项只需一个公开多项式乘数，以及自同构后的密钥切换，不涉及解密，也不除以扩域次数。特别地，当 $p\mid d$ 时，不能把槽位迹除以 $d$。提升的 Frobenius 是在生成元上按 $X\mapsto X^p$ 定义的自同构；当 $j>1$ 时，它不是对模 $p^j$ 下任意元素直接应用 $W\mapsto W^p$。

对 $0\leq u<d$，通过 CRT 剩余定义 $\mathcal T_u$：槽位 $h$ 放入标量系数 $z_{u\ell+h}$，其余扩域坐标全为零。逆映射 $\mathcal U_u$ 读取各槽位的标量坐标，将其写入系数位置 $u\ell+h$，其他位置补零。CRT 矩阵求逆即可构造这些映射，迹对偶公式则将正逆映射都转化为同态操作。它们满足
\[
\sum_{u=0}^{d-1}\mathcal U_u\bigl(\mathcal T_u(Z)\bigr)=Z.
\]
打包映射在模 $p^\varepsilon$ 下使用，相容的解包映射在模 $p^r$ 下使用。数字多项式作用于标量项，而不是任意扩域元素。若其分子的每个槽位都能被 $p$ 整除，CRT 就推出整个多项式属于 $pR_j$，因此后续逐系数精确除法有良好定义。这给出了下述变换的一般但未经优化的构造~\cite{cryptoeprint:2022/1363}。

熟悉的完全分裂矩阵是 $d=1$ 的特殊情形：

\textsf{CoeffToSlot} 步骤的目标是以同态方式将多项式 $Z$ 的系数 $z_i$ 移入输入向量槽位。


在小结~\ref*{subsubsec:bfv-rotation-summary}（位于 \autoref{subsubsec:bfv-rotation-summary}）中，我们了解到，将输入槽位向量 $\vec{v}$ 转换为多项式系数向量 $\vec{m}$ 的编码公式为 $\vec{m} = n^{-1}\cdot\hathat{W} \cdot I_n^R \cdot \vec{v}$，其中，$\hathat{W}$ 是按如下方式构造的 $n$ 维向量空间的一组基：

$\hathat{W} = \begin{bmatrix}
1 & 1 & \cdots & 1 & 1 & 1 & \cdots & 1\\
(\omega^{J(\frac{n}{2} - 1)}) & (\omega^{J(\frac{n}{2} - 2)}) & \cdots & (\omega^{J(0)}) & (\omega^{J_*(\frac{n}{2} - 1)}) & (\omega^{J_*(\frac{n}{2} - 2)}) & \cdots & (\omega^{J_*(0)})\\
(\omega^{J(\frac{n}{2} - 1)})^2 & (\omega^{J(\frac{n}{2} - 2)})^2 & \cdots & (\omega^{J(0)})^2 & (\omega^{J_*(\frac{n}{2} - 1)})^2 & (\omega^{J_*(\frac{n}{2} - 2)})^2 & \cdots & (\omega^{J_*(0)})^2 \\
\vdots & \vdots & \ddots & \vdots & \vdots & \vdots & \ddots & \vdots \\
(\omega^{J(\frac{n}{2} - 1)})^{n-1} & (\omega^{J(\frac{n}{2} - 2)})^{n-1} & \cdots & (\omega^{J(0)})^{n-1} & (\omega^{J_*(\frac{n}{2} - 1)})^{n-1} & (\omega^{J_*(\frac{n}{2} - 2)})^{n-1} & \cdots  & (\omega^{J_*(0)})^{n-1}
\end{bmatrix}$

$ $

\textcolor{red}{ $\rhd$ 其中，旋转辅助函数为 $J(h) = 5^h \bmod 2n$}

$ $

因此，对于输入密文 $\textsf{ct} = \textsf{RLWE}_{S, \sigma}\bm(\Delta' Z\bm) \bmod q$，可以将其输入向量槽位理解为存储了某些值，对它们乘以 $n^{-1}\cdot\hathat{W} \cdot I_n^R$ 后就得到多项式 $Z$ 的系数 $z_i$。这意味着，如果对 $\textsf{RLWE}_{S, \sigma}\bm(\Delta' Z\bm)$ 的输入向量槽位\textit{同态}乘以 $n^{-1}\cdot\hathat{W} \cdot I_n^R$，那么所得密文的 $n$ 个输入向量槽位将包含 $Z$ 的 $n$ 个系数，等价于将 $Z$ 的系数移入输入向量槽位。因此，\textsf{CoeffToSlot} 步骤等价于同态计算 $n^{-1}\cdot\hathat{W} \cdot I_n^R \cdot \textsf{RLWE}_{S, \sigma}\bm(\Delta' Z\bm)$。可以使用 \autoref{subsec:bfv-matrix-multiplication} 中介绍的技术同态计算矩阵与向量相乘。

在 \textsf{CoeffToSlot} 步骤之后，可以通过对去噪多项式进行同态求值，消除每个 $z_i$ 的低 $\varepsilon-r$ 个（$p$ 进制）数位中的噪声；下一小节将介绍这一多项式。

在明文槽位中恢复 $M$ 的系数后，需要将它们移回原来的系数位置。密文仍含有有界的求值误差。这一步称为 \textsf{SlotToCoeff}，它恰好是 \textsf{CoeffToSlot} 的逆过程。由小结~\ref*{subsubsec:bfv-rotation-summary}（位于 \autoref{subsubsec:bfv-rotation-summary}）还可知，$n^{-1}\cdot\hathat{W} \cdot I_n^R$ 的逆矩阵为 $\hathat{W}^*$，其中：

$\hathat{W}^* = \begin{bmatrix}
1 & (\omega^{J(0)}) & (\omega^{J(0)})^2 & \cdots & (\omega^{J(0)})^{n-1}\\
1 & (\omega^{J(1)}) & (\omega^{J(1)})^2 & \cdots & (\omega^{J(1)})^{n-1}\\
1 & (\omega^{J(2)}) & (\omega^{J(2)})^2 & \cdots & (\omega^{J(2)})^{n-1}\\
\vdots & \vdots & \vdots & \ddots & \vdots \\
1 & (\omega^{J(\frac{n}{2}-1)}) & (\omega^{J(\frac{n}{2}-1)})^2 & \cdots & (\omega^{J(\frac{n}{2}-1)})^{n-1}\\
1 & (\omega^{J_*(0)}) & (\omega^{J_*(0)})^2 & \cdots & (\omega^{J_*(0)})^{n-1}\\
1 & (\omega^{J_*(1)}) & (\omega^{J_*(1)})^2 & \cdots & (\omega^{J_*(1)})^{n-1}\\
1 & (\omega^{J_*(2)}) & (\omega^{J_*(2)})^2 & \cdots & (\omega^{J_*(2)})^{n-1}\\
\vdots & \vdots & \vdots & \ddots & \vdots \\
1 & (\omega^{J_*(\frac{n}{2}-1)}) & (\omega^{J_*(\frac{n}{2}-1)})^2 & \cdots & (\omega^{J_*(\frac{n}{2}-1)})^{n-1}\\
\end{bmatrix}$

因此，\textsf{SlotToCoeff} 步骤等价于将去噪多项式求值的输出同态乘以 $\hathat{W}^*$。

下一小节将学习如何基于数位提取多项式 $G_\varepsilon(x)$，设计 BFV 的核心算法——去噪多项式。


\subsubsection{数位提取}
\label{subsubsec:bfv-bootstrapping-digit-extraction}


当 $p$ 为奇数时，使用平衡的 $p$ 进制数字 $d\in\{-(p-1)/2,\ldots,(p-1)/2\}$。将输入系数写成 $z=\sum_{j=0}^{\varepsilon-1}d_jp^j\pmod{p^\varepsilon}$。删除最低 $v=\varepsilon-r$ 个平衡数字，得到解密所需的 $\lceil z/p^v\rfloor\pmod{p^r}$。普通非负数字实现的则是向下取整：例如 $p=3,\varepsilon=3,m=1,e'=-1$ 时 $z=8$，虽然 $|e'|<9/2$，仍有 $\lfloor8/9\rfloor=0\ne1$。当 $p=2$ 时，使用非负数字，并在删除 $v$ 位之前加上 $2^{v-1}$，从而在非中点处实现 $\lceil z/2^v\rfloor=\lfloor(z+2^{v-1})/2^v\rfloor$。这里的 $p=2$ 公式仅是标量舍入恒等式，不能把所引用的 CRT 自举路线推广到本章的二次幂环；该路线要求分圆指标与 $p$ 互素，详见 \autoref{subsubsec:bfv-bootstrapping-coefftoslot}。

对 $j\geq1$，提升多项式 $F_j$ 满足
\[
F_j(d+p^ju)\equiv d\pmod{p^{j+1}}
\]
其中 $d$ 为允许的最低位数字，$u$ 为整数。下标 $j$ 表示算法中已知的精度，不必等于秘密输入中首个非零高位的位置。令 $L_\varepsilon=F_{\varepsilon-1}\circ\cdots\circ F_1$。

\para{第 1 次递归：}$F_1(z)\equiv d_0\pmod{p^2}$。

\para{第 2 次递归：}$F_2\circ F_1(z)\equiv d_0\pmod{p^3}$。

\para{第 3 次递归：}$F_3\circ F_2\circ F_1(z)\equiv d_0\pmod{p^4}$。

\para{第 $\bm{(\varepsilon-1)}$ 次递归：}$L_\varepsilon(z)\equiv d_0\pmod{p^\varepsilon}$。

因此，数字删除操作为
\[
G_\varepsilon(z)=\frac{z-L_\varepsilon(z)}p\pmod{p^{\varepsilon-1}}.
\]
这里对可整除的代表执行精确整数除法，输出模数降低一个 $p$ 的幂次。它不是在 $\mathbb Z_{p^\varepsilon}$ 中乘以 $p$ 的逆元，因为该逆元不存在。即使存在 $p^{-1}\bmod q$，用它乘密文分量也是不同操作，可能放大噪声。因此，$G_\varepsilon$ 表示多项式求值与除法组成的过程，而不是 $\mathbb Z_{p^\varepsilon}$ 上的单个多项式。重复执行
\[
G_{r+1}\circ\cdots\circ G_{\varepsilon-1}\circ G_\varepsilon(z)
\]
即可删除 $\varepsilon-r$ 位。每次调用 $G_j$ 都使用 $F_1,\ldots,F_{j-1}$ 提取自身输入的最低位；即使 $r>1$，也必须从 $F_1$ 开始。

\para{尺度因子的重新解释：}假设同态减法后，明文代表满足 $u=pw\pmod{p^j}$，密文相位为 $\Delta_j u+E_j$，其中 $\Delta_j=\lfloor q/p^j\rfloor$。选择一致且可整除的中心代表后，同一密文可在明文模数 $p^{j-1}$ 下重新解释为
\[
\Delta_jpw+E_j=\Delta_{j-1}w+(E_j-\delta_jw),\qquad
\delta_j=\Delta_{j-1}-p\Delta_j\in\{0,\ldots,p-1\}.
\]
这避免了将密文分量乘以 $p^{-1}\bmod q$。其中 $-\delta_jw$ 是按地板尺度衡量误差时的坐标变化；若按精确尺度衡量，它恰好抵消：

\[
\frac{q}{p^j}(pw)+\mathcal E_j
=\frac{q}{p^{j-1}}w+\mathcal E_j.
\]
因此，精确除法不增加自然误差，而且扩大了解码容差。两种描述等价：若误差界已经采用自然误差坐标，就不能再次累加地板尺度的漂移。使用相容且能被整除的中心系数代表时，这个恒等式精确成立，包括 $p^j\nmid q$ 的情形。

各地板尺度依次为 $\lfloor q/p^\varepsilon\rfloor,\lfloor q/p^{\varepsilon-1}\rfloor,\ldots,\lfloor q/p^r\rfloor$。与密文相加的明文必须采用当前尺度，而密文与明文相乘使用未缩放乘数，如 \autoref{subsec:bfv-mult-plain} 所述。多项式求值会改变自然误差；精确除法保持自然误差，同时改变明文模数。\autoref{subsubsec:bfv-bootstrapping-summary} 的显式递推可逐项检查所有中间解码条件。

剩余的局部代数任务是构造提升多项式。

\para{设计 \boldmath$F_{\varepsilon'}(z_i)$：}下面给出显式控制精度的插值构造。令 $j=\varepsilon'\geq1$，$D$ 为选定的连续 $p$ 个数字代表组成的集合。由费马小定理，对每个 $d\in D$，$(d^p-d)/p$ 都是整数。在 $\mathbb Z_{p^j}$ 上插值得到次数至多为 $p-1$ 的多项式 $H_j$，使得
\[
H_j(d)\equiv(d^p-d)/p\pmod{p^j}\quad(d\in D).
\]
不同数字节点之差在模 $p^j$ 下都是单位元，因此拉格朗日插值成立。选取系数的整数提升，并定义
\[
F_j(X)=X^p-pH_j(X).
\]
\begin{proof}
对于 $z=d+p^ju$，二项式定理给出 $z^p\equiv d^p\pmod{p^{j+1}}$。整数系数多项式的求值还满足 $H_j(z)\equiv H_j(d)\pmod{p^j}$。因此
\[
F_j(z)\equiv d^p-pH_j(d)\equiv d\pmod{p^{j+1}}.
\]
\end{proof}
这证明了各次递归所用的提升性质。其次数至多为 $p$，不是 $p-1$。若仅在 $\mathbb F_p$ 上分别插值各位数字函数，再重复使用任意整数系数提升，则不足以保证更高幂次下的正确性，因为进位会影响高位。系数必须按上述方法在所需精度下构造。这一局部存在性证明本身并没有给出优化同态实现的深度界或噪声界。

\subsubsection{小结}
\label{subsubsec:bfv-bootstrapping-summary}

下面总结推广后的明文变换。方框后的递推明确给出正确性条件，涵盖低模舍入、一般线性变换、全部同态门以及最终解码。安全性与优化仍是独立问题。

\begin{tcolorbox}[title={\textbf{\tboxlabel{\ref*{subsubsec:bfv-bootstrapping-summary}} BFV 自举}}]
令 $t=p^r$、$\varepsilon>r$、$P=p^\varepsilon$、$\Delta=\lfloor q/p^r\rfloor$ 和 $\Delta'=\lfloor q/P\rfloor$。 完整 CRT 流程须满足 \autoref{subsubsec:bfv-bootstrapping-coefftoslot} 中的相容环条件；在本章的二次幂环中，$p$ 为奇数。
\begin{enumerate}
\item \textbf{模数切换：}计算 $A'=[\lceil(P/q)A\rfloor]_P$ 和 $B'=[\lceil(P/q)B\rfloor]_P$。其相位满足 $A'S+B'\equiv p^{\varepsilon-r}M+E'\pmod P$，完整 $E'$ 见 \autoref{subsubsec:bfv-bootstrapping-modulus-switch}。要求 $\|E'\|_\infty<p^{\varepsilon-r}/2$。
\item \textbf{同态解密：}使用模数 $q$ 下 $\Delta'S$ 的加密，计算 $A'\cdot\textsf{RLWE}_{S,\sigma}(\Delta'S)+(0,\Delta'B')$。按 \autoref{subsubsec:bfv-bootstrapping-homomorphic-decryption} 同时跟踪密钥噪声项和明文约减余数项。
\item \textbf{\textsf{CoeffToSlot}：}同态地将 $Z=[p^{\varepsilon-r}M+E']_P$ 的系数放入槽位。在标量分裂情形下，槽位变换为 $n^{-1}\hathat W I_n^R$（\autoref{subsubsec:bfv-rotation-summary}）；对于扩域槽位，则使用 \autoref{subsubsec:bfv-bootstrapping-coefftoslot} 构造的 $d$ 个映射 $\mathcal T_u$。
\item \textbf{数字提取：}对当前精度 $j=\varepsilon,\varepsilon-1,\ldots,r+1$，用 $L_j=F_{j-1}\circ\cdots\circ F_1$ 提取最低位，再实现精确商 $G_j(z)=(z-L_j(z))/p\pmod{p^{j-1}}$。当 $p$ 为奇数时使用平衡数字；当 $p=2$ 时使用上文规定的舍入偏移量。最终槽位为 $m_i\pmod{p^r}$，所以未约减代表可以相差 $k_i'p^r$；当 $r>1$ 时，不能写成 $k_i'p$。
\para{\underline{尺度因子的重新解释}：}每次除法都把尺度从 $\lfloor q/p^j\rfloor$ 重新解释为 $\lfloor q/p^{j-1}\rfloor$，并计入上文推导的漂移 $-\delta_jw$。最终尺度为 $\lfloor q/p^r\rfloor=\Delta$。这里不假设 $p$ 存在模逆元。
\item \textbf{\textsf{SlotToCoeff}：}在标量分裂情形下应用逆线性变换 $\hathat W^*$，扩域槽位情形则对 $\mathcal U_u$ 的输出求和，恢复相位为 $\Delta M+E^{\langle b\rangle}\pmod q$ 的加密。完整输出误差包含全部求值误差、尺度漂移及明文约减项，并须满足 BFV 解码界。
\end{enumerate}
\end{tcolorbox}

\textbf{显式的条件性正确性证书。}采用 \autoref{subsubsec:bfv-bootstrapping-high-level} 定义的自然误差 $\mathcal E_T$。令 $C_R$ 满足 $\|UV\|_\infty\leq C_R\|U\|_\infty\|V\|_\infty$，并令 $h_S$ 是乘以秘密多项式这一线性算子的无穷范数上界。对于负循环环，可取 $C_R=n$ 和 $h_S=\|S\|_1$。对于公开多项式 $C$ 或自同构 $\tau_a$，分别以 $\lambda_C$、$\lambda_a$ 表示相应系数算子的无穷范数。这些都是可计算的矩阵范数，也适用于一般分圆环。

使用具有 $L$ 个数字、数字绝对值不超过 $D_g$ 的精确整数分解，求值密钥各行在输出秘密下的原始相位为 $g_iS_{\rm old}+E_i^{\rm key}$，且 $\|E_i^{\rm key}\|_\infty\leq B_{\rm key}$。一次密钥切换至多给自然误差上界增加
\[
B_{\rm ks}=C_RLD_gB_{\rm key}
\]
这里不需要独立性假设。若输入界为 $B_1,B_2$，可递推下列有效误差界；必要时可以将界向上取整：
\[
\begin{aligned}
B_{\rm add/sub}&=B_1+B_2,\qquad B_{\rm plain\,mult}=\lambda_CB_1,\\
B_{\rm plain\,add}&=B_1+(r_T/T)\|C\|_\infty,\\
B_{\rm auto}&=\lambda_aB_1+B_{\rm ks},\qquad B_{\rm exact\,div}=B_1,\\
\kappa_i&=(h_S+2)/2+B_i/q.
\end{aligned}
\]
明文加法使用 $\lfloor q/T\rfloor C$；精确除法要求明文能被 $p$ 整除，并将 $T$ 改为 $T/p$。标准 FV 乘法给出
\[
\begin{aligned}
B_{\rm mult}={}&C_R\left[\frac T2(B_1+B_2)
+T(\kappa_1B_2+\kappa_2B_1)+\frac TqB_1B_2\right]\\
&+\frac{1+h_S+h_S^2}{2}+B_{\rm ks}.
\end{aligned}
\]
为核实此界，将中心分量提升写成 $AS+B=(q/T)M+\mathcal E+qK$，由范数可得 $\|K\|_\infty\leq\kappa$。展开整数张量乘积，再乘以 $T/q$，保留下列自然误差项：
\[
\begin{aligned}
\mathcal E_{\rm mult}={}&M_1\mathcal E_2+M_2\mathcal E_1
+T(\mathcal E_1K_2+\mathcal E_2K_1)+\frac Tq\mathcal E_1\mathcal E_2\\
&+\rho_0+\rho_1S+\rho_2S^2+E_{\rm ks}\pmod q.
\end{aligned}
\]
省略的进位项都是 $q$ 的整数倍。各张量分量舍入误差 $\rho_i$ 的范数至多为 $1/2$，从而得到上述舍入界。明文模 $T$ 约减也只使精确缩放的明文改变 $q$ 的倍数。这就证明了递推式，而不是预先假定乘法正确。

对于最初的模数切换，输入自然误差界 $B_{\rm in}$ 满足下式即可：
\[
\frac PqB_{\rm in}+\frac{1+h_{\rm in}}2
<\frac{p^{\varepsilon-r}}2.
\]
这里 $h_{\rm in}$ 是乘以输入秘密这一算子的范数上界。旧误差此时成为整数明文 $Z$ 的一部分，并不是新求值密文的加密误差。若输入秘密的加密具有自然误差界 $B_s$，则同态解密后的求值误差账本从下式开始：
\[
B_Z\leq\lambda_{A'}B_s+(r_P/P)\|B'\|_\infty.
\]
若自举密钥使用地板尺度加密，且整数误差界为 $B_{\rm key}$，可取 $B_s=B_{\rm key}+(r_P/P)\|S_{\rm in}\|_\infty$。求值模数可以独立于输入模数选取；为简化记号，上式使用相同的 $q$。

要求显式电路中每个密文节点的界均满足 $B_C<q/(2T_C)$。模数切换不等式保证选择正确的舍入明文；迹对偶映射将其系数放入标量槽位；提升多项式和精确除法恢复模 $p^r$ 的各系数；逆映射重构 $M$。结合上述递推归纳，即可证明最终密文解密为 $M$。在地板尺度记号下，其输出误差为 $E^{\langle b\rangle}=\mathcal E_{p^r}+(r_{p^r}/p^r)M$。新误差界不必小于每一种输入误差，但必须满足后续计算的预算。

\texttt{audit/check\_bootstrap.py} 中的精确公开求值样例覆盖了完全标量分裂的 $\Phi_4,p=5$、在 $r=1$ 和 $r=2$ 下具有二次扩域槽位的 $\Phi_8,p=3$，以及互素二进制替代环 $\Phi_3,p=2$。各例的密文模数均不能被明文模数整除，输入误差和求值密钥误差也都非零。独立检查器在公开求值完成后审计每个中间密文；求值器不接收秘密或解密回调。这些小环样例验证了实现与代数，并不证明密码学安全性，也没有实现参考文献中优化的薄自举算法。

\para{同态解密的目的：}它将低模数密文中依赖秘密密钥的相位，转换为新求值模数及适当尺度下的加密明文，使后续舍入电路能够在不泄露秘密密钥和明文的情况下求值。尺度重新解释只是数字删除实现的一个组成部分，本身并不构成自举的完整解释或正确性证明。

